\section{Conditional route: dual lattice current realization}
\label{app:dlcr_conditional}

This appendix records a new candidate mechanism for the two live realization
theorems. It is not an established theorem package in the present repository.
It should be read as a precise research target: if the lemmas below are proved
with the stated compatibility, then the two packet realization blockers in
Appendix~\ref{sec:hamiltonian-identification-boundary} would be discharged.

\subsection{Finite current charts}
For each finite Wilson-loop packet $W_N$ and rational-time packet $F_N$, the
proposed construction introduces a finite dual-current chart $J_N$. The intended
role of $J_N$ is to carry both Wilson-loop packet observables and rational-time
time-translation representatives in one finite-dimensional gauge-invariant
observable space. The first required lemma is:
\begin{quote}
For every $N$, each packet observable factors through $J_N$ up to a gauge-exact
cylindrical coboundary.
\end{quote}
This would turn both realization problems into finite current-law comparison
problems.

A finite-coordinate version of this step is discharged first. Since $W_N$ and
$F_N$ are finite, enumerate all packet observables as $O_1,\ldots,O_m$ and set
\[
  J_N(A)=(O_1(A),\ldots,O_m(A)).
\]
Every packet observable is exactly recovered by a coordinate projection. This
proves finite-dimensional packet packaging.

The finite gauge-theoretic upgrade is also discharged on each packet graph.
Let $\Gamma_N$ be the finite graph spanned by the packet paths and choose a
maximal tree $T$. Tree gauge sets every tree-edge holonomy to the identity.
The remaining non-tree edge variables form a finite cycle-current chart with
\[
  |E(\Gamma_N)|-|V(\Gamma_N)|+c(\Gamma_N)
\]
group-valued coordinates. Every packet path holonomy is a word in these
non-tree variables and their inverses. A closed Wilson loop changes only by
residual root conjugation, and its normalized trace is therefore unchanged. The
factorization error is exactly zero. Thus the finite DLCR chart lemma is proved
packetwise; the remaining DLCR work is analytic, not finite graph packaging.

\subsection{Exact shadow coupling}
The second required lemma is an exact shadow coupling. Its finite probability
core is discharged. If the convergence-side chart law and the OS chart law on
$J_N$ are equal, the diagonal coupling gives equality almost surely and hence
zero error for all packet tests. More generally, if their total variation
distance is $\delta_N(a)$, then every bounded packet test $\Phi$ obeys
\[
  \left|\mathbf E_{\mathrm{conv}}\Phi-\mathbf E_{\mathrm{OS}}\Phi\right|
  \leq 2\|\Phi\|_\infty\delta_N(a).
\]

Thus the remaining shadow-coupling problem is a chart-law theorem for
Yang--Mills. Equivalently, it is the finite chart-law convergence problem.
This problem has a standard compact-group reduction. By Peter--Weyl, the
gauge-invariant matrix-coefficient polynomials are uniformly dense in the
continuous class functions on the finite chart $SU(3)^m/\!\operatorname{Ad}$.
If a packet test $\Phi$ is approximated by a finite Peter--Weyl polynomial $P$
with $\|\Phi-P\|_\infty\leq\varepsilon$, then
\[
  \left|\mathbf E_\mu\Phi-\mathbf E_\nu\Phi\right|
  \leq \left|\mathbf E_\mu P-\mathbf E_\nu P\right|+2\varepsilon.
\]
Thus finite chart-law convergence reduces to convergence of finitely truncated
spin-network moments plus a uniform approximation tail.

The finite chart-law tests can be made more concrete on the same packet graph.
After tree gauge fixing, let $x_1,\ldots,x_m$ denote the non-tree
$SU(3)$ variables and, for a word $w$ in the free group on these variables, set
\[
  \tau_w(x)=\frac{1}{3}\operatorname{Tr}(w(x_1,\ldots,x_m)).
\]
The unital $*$-algebra generated by the $\tau_w$ is dense in
$C(SU(3)^m/\!\operatorname{Ad})$. Indeed it is closed under conjugation since
$\overline{\tau_w}=\tau_{w^{-1}}$, and it separates quotient points: equality
of all trace words is equality of the characters of the corresponding
three-dimensional unitary representations of the free group. Finite-dimensional
unitary representations are semisimple and irreducible characters are linearly
independent, hence the representations are unitarily equivalent, and the
conjugating unitary can be rescaled to determinant one.
Stone--Weierstrass applies. Therefore, if
$Q=\sum_\ell c_\ell M_\ell$ is a finite polynomial in trace words and
$\|\Phi-Q\|_\infty\leq\varepsilon$, then
\[
  \left|\mathbf E_\mu\Phi-\mathbf E_\nu\Phi\right|
  \leq
  \sum_\ell |c_\ell|
  \left|\mathbf E_\mu M_\ell-\mathbf E_\nu M_\ell\right|
  +2\varepsilon .
\]
Thus the chart-law problem is now reduced to finite mixed Wilson trace-word
moment convergence plus the explicit approximation tail.

The resulting finite certificate has explicit constants. Allow the approximant
to depend on the lattice spacing,
$Q_a=\sum_\ell c_\ell(a)M_\ell(a)$; a fixed polynomial is the special case. If
the selected monomials satisfy
\[
  \left|\mathbf E_a M_\ell(a)-\mathbf E_{\mathrm{OS}}M_\ell(a)\right|
  \leq B_\ell(a)a^{\beta_\ell(a)},\qquad
  \beta_\ell(a)\geq\beta>0,
\]
with uniform coefficient-weighted constant
$\sum_\ell |c_\ell(a)|B_\ell(a)\leq B$, and the approximation tail is
power-controlled, $\|\Phi-Q_a\|_\infty\leq E a^\gamma$, then
\[
  \left|\mathbf E_a\Phi-\mathbf E_{\mathrm{OS}}\Phi\right|
  \leq
  \left(B+2E\right)a^\alpha,
  \qquad
  \alpha=\min(\beta,\gamma).
\]
If only $\|\Phi-Q_a\|_\infty\leq\varepsilon$ is available, the same argument gives
the weaker bound
\[
  \left|\mathbf E_a\Phi-\mathbf E_{\mathrm{OS}}\Phi\right|
  \leq
  Ba^\beta
  +2\varepsilon .
\]
Thus a $C_Na^{\alpha_N}$ chart-law estimate now follows from finite
trace-word moment estimates and a power-controlled trace-polynomial
approximation scheme. Those two Yang--Mills analytic inputs remain proved in
this repository.

There is one further finite reduction for the moment estimates. Put
\[
  d_\infty(x,y)=\max_j\|x_j-y_j\|_{\operatorname{op}}
\]
on the finite chart representatives. A normalized trace word
$\tau_w=3^{-1}\operatorname{Tr}(w)$ of word length $|w|$ is $|w|$-Lipschitz:
the product telescoping identity changes one letter at a time, inverses
preserve operator-norm distance on unitary matrices, and
$|3^{-1}\operatorname{Tr}A|\leq\|A\|_{\operatorname{op}}$. Hence a monomial
$M=\prod_r\tau_{w_r}$ is Lipschitz with constant at most
$\sum_r |w_r|$. Consequently any finite-chart coupling with
\[
  \mathbf E\,d_\infty(X_a,Y)\leq D a^\eta
\]
implies
\[
  \left|\mathbf E M(X_a)-\mathbf E M(Y)\right|
  \leq \left(\sum_r |w_r|\right)D a^\eta .
\]
Thus the trace-word moment problem may equivalently be attacked by constructing
a finite-chart transport coupling. The repository does not currently construct
that Yang--Mills coupling.

For the original finite Wilson-loop packet moments, the approximation part is
even sharper. Each packet loop is already a single trace word after tree gauge,
and a finite mixed packet moment is already a monomial
$M=\prod_r\tau_{w_r}$. Thus in the trace-polynomial certificate one may take
$Q=M$ and the approximation tail is exactly zero:
\[
  \|M-Q\|_\infty=0.
\]
Consequently, for packet moments, the remaining chart-law input is precisely
the finite-chart transport estimate (or an equivalent direct moment estimate);
there is no separate trace-polynomial approximation tail to prove for those
tests.

This gives a sharp finite criterion for the state-side packet realization. Let
$X_a$ be the convergence-side chart variable and $Y$ the OS chart variable. If
for the fixed packet chart there is a coupling with
\[
  \mathbf E\,d_\infty(X_a,Y)\leq D_Na^{\eta_N},\qquad 0<a\leq1,
\]
then every mixed packet moment $M=\prod_r\tau_{w_r}$ satisfies
\[
  \left|\mathbf E_aM-\mathbf E_{\mathrm{OS}}M\right|
  \leq
  L(M)D_Na^{\eta_N},
  \qquad L(M)=\sum_r |w_r|.
\]
Since $W_N$ is finite, the whole packet has the uniform constant
\[
  C_N=D_N\max_{M\in W_N}L(M),\qquad \alpha_N=\eta_N.
\]
Thus the finite packet Schwinger-realization blocker is now equivalent, in this
route, to proving the finite-chart Yang--Mills transport estimate above.

For every fixed $N$, one must prove that the lattice Yang--Mills field and the
continuum Schwinger/OS field have $J_N$-marginals satisfying an explicit estimate
\[
  \left|\mathbf E_a\Phi(J_N)-\mathbf E_{\mathrm{OS}}\Phi(J_N)\right|
  \leq C_N a^{\alpha_N}
\]
or equivalently a comparable chart-law convergence bound. Passing $a\to0$ would
then identify the convergence-side packet law with the OS Schwinger law on
$W_N$.

\subsection{Reflection-time intertwiner}
The action-side blocker requires a separate compatibility statement. Its finite algebraic core is now discharged. Enlarge $\Gamma_N$ so it contains every
representative path $d_j$ and every translated representative
$\alpha_{q_i}(d_j)$ appearing in $F_N$. In the same maximal-tree gauge, both
$d_j$ and $\alpha_{q_i}(d_j)$ are words in the non-tree cycle variables. Define
the finite packet map $A_{q_i,N}$ by word substitution,
\[
  \operatorname{word}(d_j)\mapsto \operatorname{word}(\alpha_{q_i}(d_j)).
\]
Then, entrywise on $F_N$,
\[
  J_N(\alpha_{q_i}(d_j))=A_{q_i,N}(J_N(d_j)),
\]
again with zero finite-chart error after passing to Wilson-loop traces.

The remaining action-side blocker is analytic rather than finite algebraic. The needed
intertwiner is
\[
  J_N(T_a(q)O)\longrightarrow \alpha_q(J_N(O))
\]
for every rational time $q$ appearing in $F_N$, with convergence in the same
finite chart used by the shadow coupling. If proved, this would identify the
convergence-side rational-time action with the OS Euclidean time action on the
packet.

The finite transport criterion has the same form as on the state side. On the
enlarged chart, each translated representative is an exact word, so each
translated packet test $M_q=\prod_r\tau_{w_{q,r}}$ has Lipschitz constant
$L(M_q)=\sum_r |w_{q,r}|$. If the convergence-side translated chart variable
$X_{a,q}$ and the OS translated chart variable $Y_q$ admit a coupling with
\[
  \mathbf E\,d_\infty(X_{a,q},Y_q)\leq D_{N,q}a^{\eta_{N,q}},
\]
then
\[
  \left|\mathbf E_aM_q-\mathbf E_{\mathrm{OS}}M_q\right|
  \leq L(M_q)D_{N,q}a^{\eta_{N,q}}.
\]
For the finite packet, take
\[
  C_N=\max_{q,M_q\in F_N}L(M_q)D_{N,q},
  \qquad
  \alpha_N=\min_{q,M_q\in F_N}\eta_{N,q}.
\]
Thus the action-side finite chart problem is reduced to constructing these
translated couplings and matching them with the OS time action.

\subsection{Multiscale transport-flow collapse}
There is one additional invented reduction layer that may be useful if the
bridge influence is naturally organized by scale rather than by one-shell
recurrence. Suppose the Ward-projected bridge operator is decomposed into a
finite family of scale operators $T_0,\ldots,T_{S-1}$ on the same block
partition. If the weighted norms satisfy $\prod_s \theta_s<1$, with
\[
  \theta_s=\|T_s\|_\infty,
\]
then the renormalized operator
\[
  T_{\mathrm{ren}}=T_{S-1}\cdots T_0
\]
is contractive and the screened block norm obeys
\[
  \sup_x e^{m k_N(x)} b_x \leq C a^\beta/(1-\prod_s\theta_s).
\]
This is not proved for Yang--Mills in the repository. It is a proposed
multiscale route: a transport-flow collapse principle, if established, would
imply the block influence bound, which then feeds the existing screened-norm
and leakage certificates.

The contraction test can be made still more local in scale. Suppose one proves
logarithmic scale budgets
\[
  \log\theta_s\leq \kappa_s .
\]
Then
\[
  \prod_s\theta_s\leq \exp\Bigl(\sum_s\kappa_s\Bigr).
\]
Therefore a strictly negative total budget
\[
  K:=\sum_s\kappa_s<0
\]
is enough to give
\[
  \sup_x e^{m k_N(x)} b_x
  \leq
  \frac{C}{1-e^K}a^\beta .
\]
This log-budget version is useful because some scales may expand
($\kappa_s>0$), provided other Ward-projected scales contract enough to make
the total budget negative. The repository proves only this deterministic
allocation certificate; the Yang--Mills task is still to derive the actual
scale budgets uniformly in the continuum and volume limits.

A useful corollary is a sparse bad-scale test. Suppose $G$ good scales obey
\[
  \log\theta_s\leq-\gamma
\]
and $B$ bad scales obey
\[
  \log\theta_s\leq b .
\]
Then the total budget is bounded by
\[
  K\leq Bb-G\gamma .
\]
Hence $Bb-G\gamma<0$ gives
\[
  \sup_x e^{m k_N(x)} b_x
  \leq
  \frac{C}{1-\exp(Bb-G\gamma)}a^\beta .
\]
When $b>0$, the equivalent sparsity condition is $B/G<\gamma/b$. This is
another deterministic certificate only: the analytic work would be to classify
the actual Yang--Mills Ward-projected scales into good and bad scales and to
prove that the bad scales are sparse enough uniformly in the cutoff limits.

The sparsity criterion can also be localized in scale. Suppose every
consecutive window of $W$ scales contains at most $M$ bad scales. Partition a
finite scale interval of length $S$ into $q=\lfloor S/W\rfloor$ full windows
and a remainder of length $r=S-qW$. Each full window contributes at most
\[
  Mb-(W-M)\gamma
\]
to the log budget, while the remainder is bounded pessimistically by $rb$.
Thus
\[
  K\leq q\bigl(Mb-(W-M)\gamma\bigr)+rb.
\]
If this upper bound is negative, then
\[
  \sup_x e^{m k_N(x)} b_x
  \leq
  \frac{C}{1-\exp(K)}a^\beta .
\]
Equivalently, the full-window gain
$\Delta=(W-M)\gamma-Mb$ must dominate the leftover expansion $rb$ through
$q\Delta>rb$. This prevents all bad scales from clustering in one place; it is
still only a deterministic reduction until the Yang--Mills bridge analysis
proves the window density estimate.

There is also a phase-robust finite-profile version. Suppose a renormalized
scale analysis bounds the logarithmic scale norms by a finite cyclic profile
\[
  \kappa_0,\ldots,\kappa_{P-1}
\]
with negative period budget
\[
  K_P=\sum_{i=0}^{P-1}\kappa_i<0 .
\]
For $S=qP+r$ scales, define the worst cyclic remainder
\[
  R_r=\max_j\sum_{i=0}^{r-1}\kappa_{j+i\ {\rm mod}\ P}.
\]
Then the total log budget is bounded by
\[
  K\leq qK_P+R_r .
\]
If $qK_P+R_r<0$, the screened norm obeys
\[
  \sup_x e^{m k_N(x)} b_x
  \leq
  \frac{C}{1-\exp(qK_P+R_r)}a^\beta .
\]
This certificate is insensitive to the cutoff phase. The hard Yang--Mills
content would be proving such a finite Ward-projected scale profile with
negative period budget and uniform worst-remainder control.

A coarser route would be a negative average budget with a boundary defect.
Suppose the total scale log budget for $S$ scales satisfies
\[
  K_S\leq A-\gamma S,\qquad A\geq0,\quad \gamma>0 .
\]
If $A-\gamma S<0$, then
\[
  \sup_x e^{m k_N(x)} b_x
  \leq
  \frac{C}{1-\exp(A-\gamma S)}a^\beta .
\]
With a safety margin $\delta>0$, it is enough that
$S>(A+\delta)/\gamma$. This would allow a future subadditive or
renormalization argument to prove contraction without identifying a periodic
profile; the missing analytic content is the actual Yang--Mills proof of the
negative average estimate.

A more structural way to prove such a negative average is a telescoping
Lyapunov potential. If there are scale potentials $\Phi_s$ such that
\[
  \log\theta_s\leq \Phi_s-\Phi_{s+1}-\gamma,\qquad \gamma>0,
\]
then summing over $s=0,\ldots,S-1$ gives
\[
  K_S\leq \Phi_0-\Phi_S-\gamma S .
\]
Consequently, if $\Phi_0-\Phi_S-\gamma S<0$, then
\[
  \sup_x e^{m k_N(x)} b_x
  \leq
  \frac{C}{1-\exp(\Phi_0-\Phi_S-\gamma S)}a^\beta .
\]
This would let a future proof trade local scale expansion against decreases
of a Ward-projected energy potential. The repository proves only the
telescoping implication, not the Yang--Mills potential.

The Lyapunov check can itself be made finite. Suppose the possible
Ward-projected scale regimes form a finite directed graph and each allowed
transition $i\to j$ has log-norm bound $\log\theta_{ij}$. If there are
potentials $\Phi_i$ on the finite state set such that
\[
  \log\theta_{ij}\leq \Phi_i-\Phi_j-\gamma
\]
for every allowed transition, then every admissible scale path of length $S$
satisfies
\[
  K_S\leq \operatorname{osc}(\Phi)-\gamma S .
\]
Thus $\operatorname{osc}(\Phi)-\gamma S<0$ implies
\[
  \sup_x e^{m k_N(x)}b_x
  \leq
  \frac{C}{1-\exp(\operatorname{osc}(\Phi)-\gamma S)}a^\beta .
\]
This is a finite Bellman/Lyapunov certificate. The hard analytic step is to
code the actual Yang--Mills bridge scales by finitely many regimes and verify
the transition inequalities.

Equivalently, one may check negative cycle means on the same finite directed
graph. If every directed cycle $C$ satisfies
\[
  \sum_{e\in C} w_e\leq-\gamma |C|,
\]
then any scale path of length $S$ decomposes into erased cycles plus a simple
residual path. The residual has at most $|V|-1$ edges, so with
$A\leq(|V|-1)\max_e w_e^+$ one obtains
\[
  K_S\leq A-\gamma S .
\]
Thus, for $A-\gamma S<0$, the negative-average certificate gives the same
screened norm bound. This cycle-mean formulation turns the finite-state
Lyapunov search into finitely many graph inequalities.

The converse finite-graph construction is also deterministic. Shift edge
weights by $\gamma$, writing $u_e=w_e+\gamma$. If every directed cycle has
nonpositive $u$-weight, define $d_j$ to be the maximum $u$-weight of a path
from a zero-weight super-source to $j$. The maximum is finite, and every edge
$i\to j$ satisfies $d_j\geq d_i+u_{ij}$. Setting $\Phi_j=-d_j$ gives
\[
  w_{ij}\leq \Phi_i-\Phi_j-\gamma .
\]
Thus negative cycle means are equivalent, at the finite graph level, to the
existence of a Bellman/Lyapunov potential. The remaining Yang--Mills work is
to produce the finite scale-state graph and prove its cycle inequalities.

The cycle inequalities can be checked by a finite maximum-cycle-mean
certificate. Let $D_k(v)$ be the maximum weight of a length-$k$ path ending at
$v$, with $D_0(v)=0$. Karp's formula gives
\[
  \lambda_{\max}
  =
  \max_v \min_{0\leq k<n}
  \frac{D_n(v)-D_k(v)}{n-k}.
\]
If $\lambda_{\max}\leq-\gamma$, then every directed cycle has mean at most
$-\gamma$, so the cycle-mean and Lyapunov certificates apply. This reduces the
finite graph part of the contraction proof to a dynamic-programming table.

Combining the finite graph certificates gives a direct screened-norm package.
If the finite Ward-projected scale graph has edge log weights $w_e$, Karp
certifies $\lambda_{\max}\leq-\gamma$, and paths have length $S$ with transient
defect $A\leq(|V|-1)\max_e w_e^+$ satisfying $A-\gamma S<0$, then
\[
  C_{\mathrm{SB}}
  \leq
  \frac{C}{1-\exp(A-\gamma S)} .
\]
Thus a proved Yang--Mills finite scale graph with negative maximum cycle mean
would feed the screened bridge norm route directly.

The same finite graph package is robust under interval estimates. If each
edge log norm is known only through $l_e\leq w_e\leq u_e$, path weights and
cycle means are monotone in the edge weights. It is therefore enough to run
the maximum-cycle-mean and screened-norm certificate on the upper-endpoint
graph $w_e=u_e$. If that graph contracts, every graph in the certified
interval box contracts with the same screened norm constant.

Before a coarse graph is used, duplicate coarse edge declarations must be made
canonical. If several certified upper bounds $U_1,\ldots,U_r$ are declared for
the same ordered coarse edge $A\to B$, set
\[
  U_*(A,B):=\min_{1\leq q\leq r} U_q .
\]
Since the true edge log norm is bounded by each $U_q$, it is also bounded by
$U_*(A,B)$. Thus duplicate declarations may be replaced by a single
input-order-independent upper-envelope edge, and the Karp test is run on this
canonical graph. This is a certificate-management lemma only; the Yang--Mills
work remains to supply the certified edge bounds themselves.

There is also a coarse-grained version for many raw scale regimes. Let
$\pi$ map raw regimes to finitely many coarse regimes, and suppose every raw
edge $i\to j$ satisfies the upper-envelope estimate
\[
  w_{\rm raw}(i,j)\leq W_{\rm coarse}(\pi(i),\pi(j)).
\]
Then each raw path maps edge-by-edge to a coarse path, and summing the
envelope inequalities gives
\[
  K_{\rm raw}({\rm path})
  \leq
  K_{\rm coarse}(\pi({\rm path})) .
\]
Consequently a contracting coarse graph, certified for example by the Karp
maximum-cycle-mean package above, gives the same screened-norm constant for
all raw paths dominated by the envelope. This is only a deterministic
coarse-graining reduction; the Yang--Mills input would be the actual
construction of the coarse regimes, the raw-to-coarse edge envelopes, and the
negative coarse maximum cycle mean.

The raw-to-coarse envelopes can themselves be reduced to endpoint checks when
the raw edge estimates are organized by parameter boxes. Suppose on a
rectangular box $a_r\leq x_r\leq b_r$ an edge log norm is bounded by an affine
upper model
\[
  w_e(x)\leq c_e+\sum_r\alpha_{e,r}x_r .
\]
Then the maximum of the right side is attained at the signed endpoint vertex,
so
\[
  w_e(x)\leq
  U_e
  :=
  c_e
  +\sum_{\alpha_{e,r}\geq0}\alpha_{e,r}b_r
  +\sum_{\alpha_{e,r}<0}\alpha_{e,r}a_r .
\]
If these finitely many endpoint values satisfy
\[
  U_e\leq W_{\rm coarse}(\pi(i_e),\pi(j_e)),
\]
then every raw graph in the parameter boxes is edgewise dominated by the
coarse envelope, and the previous coarse-graph certificate applies. This
turns a continuum family of raw estimates into a finite endpoint computation,
but the Yang--Mills work remains to derive the affine upper models and prove
the contracting coarse graph uniformly.

A still closer finite certificate starts from nonlinear edge estimates with
derivative bounds. Let $f_e(x)$ be the raw edge log-norm bound on the same box,
let $x_0$ be a base point in the box, and suppose the Hessian is bounded in
the $\ell^1$ operator sense by $H_e$. The Taylor data must use one common
finite coordinate set:
\[
  {\rm dom}(x_0)={\rm dom}(\nabla f_e)={\rm dom}\{[a_r,b_r]\}.
\]
The base point must also satisfy $x_{0,r}\in[a_r,b_r]$ for every coordinate,
and the Hessian and any explicit remainder budget must be nonnegative. These
coordinate checks prevent a chart from covering a parameter direction that the
Taylor polynomial silently ignores, or from adding a gradient coordinate that
has no declared interval. In addition, the finite certificate uses the Hessian
remainder floor unless a separate sharper-remainder theorem is supplied: a
declared explicit remainder may not be smaller than
\[
  E_H=\frac12 H_e R_e^2 .
\]
This prevents a raw-edge witness from tightening its edge envelope by merely
declaring a smaller Taylor error term. With
\[
  R_e=\sum_r\max\{|a_r-x_{0,r}|,\ |b_r-x_{0,r}|\},
\]
Taylor's theorem gives
\[
  f_e(x)
  \leq
  f_e(x_0)+\nabla f_e(x_0)\cdot(x-x_0)
  +\frac12 H_e R_e^2
\]
throughout the box. The right side is an affine upper model, so the
box-envelope certificate applies to its signed endpoint values. Hence
uniform derivative and Hessian bounds for the raw Ward-projected edge norms,
together with a contracting coarse envelope, would give the same screened-norm
constant for every raw parameter realization. The deterministic theorem is
proved here; the Yang--Mills content is still the actual derivative/Hessian
estimate and the negative coarse cycle-mean bound.

The same argument localizes over a finite Taylor atlas. If the edge parameter
domain is covered by finitely many subboxes $B_{e,c}$ and each subbox has its
own Taylor/Hessian enclosure, the preceding endpoint calculation gives a
chart bound $U_{e,c}$. Since every admissible parameter lies in at least one
chart,
\[
  w_e(x)\leq U_e:=\max_c U_{e,c}.
\]
Thus a finite family of local Taylor estimates gives a single certified edge
upper bound. If these atlas bounds satisfy the coarse envelope inequalities,
the coarse graph again controls all raw paths. This permits highly nonlinear
raw estimates to be handled by finite subdivision, while leaving the true
Yang--Mills task as proving the chart cover, the local derivative bounds, and
the negative coarse maximum cycle mean.

For rectangular parameter domains the finite atlas cover is itself
certifiable. If each coordinate interval $[a_r,b_r]$ is partitioned by a
finite grid
\[
  a_r=t_{r,0}\leq t_{r,1}\leq\cdots\leq t_{r,n_r}=b_r,
\]
then the product cells
\[
  \prod_r [t_{r,k_r},t_{r,k_r+1}]
\]
form a finite closed cover of $\prod_r[a_r,b_r]$, with
$\prod_r n_r$ cells. Indeed, each coordinate of any point in the box lies in
one adjacent grid interval, and the product of those intervals contains the
point. Thus a Taylor atlas can be generated from one-dimensional endpoint
data. This only proves coverage; the Yang--Mills estimates on each cell remain
the analytic input.

More generally, one need not start from a product grid. Given finitely many
rectangular Taylor charts, collect every chart endpoint in each coordinate and
refine the full parameter box into elementary product cells. If every such
cell is contained in at least one chart, then the charts cover the whole
parameter domain. This endpoint-refined cover test prevents a raw-edge witness
from leaving unverified holes in parameter space; the analytic Taylor bounds on
the covering charts remain separate Yang--Mills inputs.

There is also a finite termination theorem for this atlas search. Suppose the
true raw Ward-projected edge functions are $C^2$ on a compact rectangular
parameter domain of $\ell^1$ diameter $D$, with a uniform $\ell^1$ Hessian
bound $H$. Suppose further that the proposed coarse graph has strict
raw-to-coarse slack $s$ and strict Karp margin $m_{\rm cyc}$ after reserving
the desired raw and coarse error budgets:
\[
  s_*:=s-\varepsilon_{\rm raw}-\varepsilon_{\rm coarse}>0,
  \qquad
  m_*:=m_{\rm cyc}-\varepsilon_{\rm coarse}>0 .
\]
On any Taylor chart of $\ell^1$ radius $R$, the Hessian floor is
$\frac12 H R^2$. Hence every chart radius satisfying
\[
  R\leq \sqrt{2s_*/H}
\]
when $H>0$ leaves the Taylor remainder inside the remaining slack; when
$H=0$, one chart per edge has zero Taylor remainder. Since compact rectangular
boxes admit finite product grids of arbitrarily small radius, a sufficiently
fine grid produces a finite Taylor-atlas certificate. A crude bound is
$E N^d$ charts for $E$ raw edges and $d$ parameters, with
$N\geq D/\sqrt{2s_*/H}$ in the case $H>0$. Thus, once strict margins and
uniform Hessian bounds are proved for the actual Yang--Mills raw edge
functions, the remaining certificate search is guaranteed to terminate
finitely rather than continue indefinitely. The theorem does not prove those
strict analytic margins; it proves that they are enough.

Strict margins also remove dependence on informal real-number data. Suppose a
finite real-valued Taylor-atlas/Karp certificate has raw-to-coarse slack
$s>0$ and cycle margin $m>0$, and suppose at most $M$ scalar roundoff errors
can accumulate in any checked inequality. If every scalar certificate entry is
rounded outward to a rational with denominator $Q$, then each scalar changes
by at most $1/Q$ and every inequality worsens by at most $M/Q$. Hence, whenever
\[
  \frac{M}{Q}<\min\{s,m\},
\]
the rounded rational certificate still has positive slack and positive cycle
margin. Thus any strict real certificate can be converted into exact rational
verifier data by taking $Q>M/\min\{s,m\}$. This is a credibility point for the
finite route: once the analytic inequalities are strict, the final certificate
can be made discrete, reproducible, and independent of floating-point trust.

The same rounding argument is uniform in a cutoff family. If every cutoff
certificate has raw-to-coarse slack at least $s_0>0$, cycle margin at least
$m_0>0$, at most $M$ accumulated scalar roundoff errors per checked
inequality, and common screened constants $C_{\rm SB}$ and $m$, then one
denominator
\[
  Q>M/\min\{s_0,m_0\}
\]
rounds all cutoff certificates to rational data with margins at least
$s_0-M/Q$ and $m_0-M/Q$. The screened-norm-to-leakage and
screened-norm-to-effective-action certificates then give cutoff-independent
DLCR constants. Thus cutoff-uniform strict margins would produce a uniformly
machine-checkable rational certificate family, not merely isolated finite
witnesses.

One can also certify cutoff intervals from finitely many cutoff nodes. Suppose
the cutoff range is covered by balls of radius $\rho$ around verified node
certificates, each with slack at least $s>0$ and cycle margin at least $m>0$.
If every scalar entering the certificate varies with cutoff at Lipschitz rate
at most $L$, then inside such a ball all checked inequalities can worsen by at
most
\[
  \Delta=L\rho .
\]
When $\Delta<\min\{s,m\}$, the whole ball remains certified with slack at
least $s-\Delta$ and cycle margin at least $m-\Delta$. Thus finitely many
verified cutoff nodes can certify an entire cutoff interval, provided the
cutoff Lipschitz estimates fit inside the strict margins.

Combining the preceding strict-margin reductions gives a compact analytic seed
criterion. Suppose the actual Yang--Mills raw Ward-projected edge functions are
defined on compact parameter boxes, have uniform $C^2$ bounds with Hessian
constant $H$, and admit strict raw-to-coarse slack $s>0$ and strict Karp cycle
margin $m>0$. Reserve raw and coarse error budgets
$\varepsilon_{\rm raw},\varepsilon_{\rm coarse}$ and write
\[
  s_*:=s-\varepsilon_{\rm raw}-\varepsilon_{\rm coarse},
  \qquad
  m_*:=m-\varepsilon_{\rm coarse}.
\]
Assume $s_*,m_*>0$. If cutoff variation on balls of radius $\rho$ is bounded
by $\Delta=L\rho$, and if at most $M$ scalar roundoff errors accumulate in
each checked inequality, then the single condition
\[
  s_*>\Delta+M/Q,\qquad m_*>\Delta+M/Q
\]
promotes the analytic seed to a finite, exact, cutoff-uniform certificate
program. Indeed, the Taylor-atlas termination theorem gives a finite chart
mesh with $R\leq \sqrt{2s_*/H}$ when $H>0$; the cutoff stability theorem
extends verified nodes to the covered cutoff intervals with margins
$s_*-\Delta$ and $m_*-\Delta$; and strict-margin rationalization rounds the
finite witnesses outward to denominator $Q$ while preserving exact positive
margins $s_*-\Delta-M/Q$ and $m_*-\Delta-M/Q$. Common screened constants then
feed the deterministic DLCR output certificates uniformly. This is the current
most compact remaining target: prove this strict analytic seed for the genuine
Yang--Mills raw edges. The reduction does not construct the seed, but it shows
that such a seed is sufficient for a finite reproducible certificate rather
than an open-ended search.

There is also a finite validation principle for approximate analytic evidence.
Suppose a proposed Yang--Mills seed is accompanied by rigorous interval radii.
For every upper-bound quantity, such as a raw Lipschitz constant, Hessian
constant, cutoff Lipschitz constant, or screened norm constant, replace the
nominal value $\widehat U$ by
\[
  U_{\rm cert}:=\widehat U+\operatorname{err}(U).
\]
For every lower-bound quantity, such as raw-to-coarse slack, Karp cycle margin,
or screening rate, replace the nominal value $\widehat \ell$ by
\[
  \ell_{\rm cert}:=\widehat \ell-\operatorname{err}(\ell).
\]
If the certified slack and cycle margin are positive and satisfy the strict
program inequalities
\[
  s_{\rm cert}-\varepsilon_{\rm raw}-\varepsilon_{\rm coarse}
    >(L_{\rm cert}\rho)+M/Q,
  \qquad
  m_{\rm cert}-\varepsilon_{\rm coarse}
    >(L_{\rm cert}\rho)+M/Q,
\]
then the interval evidence validates the strict analytic seed and hence
promotes to the same finite exact cutoff-uniform certificate program. This
puts a concrete target on the remaining computational/analytic work: produce
rigorous intervals with strict surplus. The validation theorem does not create
those intervals, but it proves that sufficiently sharp interval/Taylor
estimates would be accepted without any floating-point trust.

The interval validation search itself is finite under a standard convergence
model. Let $\eta>0$ be the true strict surplus in the seed program after all
reserved losses have been subtracted. Suppose a rigorous interval/Taylor
refinement scheme has aggregate validation error bounded on mesh radius $h$ by
\[
  \operatorname{Err}(h)\leq A h^p,\qquad A\geq 0,\quad p>0 .
\]
If $A=0$, any finite covering mesh suffices. If $A>0$, every mesh satisfying
\[
  h\leq \left({\eta\over 2A}\right)^{1/p}
\]
leaves certified surplus at least $\eta/2$, so the interval analytic seed
validation theorem applies. On a compact rectangular box of $\ell^1$ diameter
$D$ and dimension $d$, it is enough to take
\[
  N\geq D(2A/\eta)^{1/p}
\]
subdivisions per coordinate, giving at most $E N^d$ raw-edge charts for $E$
raw edges. Thus, once the true strict surplus and a convergent interval scheme
are established, the validation search cannot continue forever: there is an
explicit finite refinement threshold.

The strict surplus $\eta$ can itself be certified by finite residuals. Suppose
each raw edge has a certified upper bound $U_{\rm raw}(e)$ and a coarse
envelope $W_{\rm coarse}(e)$ with minimum residual
\[
  s_{\min}:=\min_e\{W_{\rm coarse}(e)-U_{\rm raw}(e)\}>0.
\]
Suppose also that the coarse graph has a finite Lyapunov potential $\Phi$ and
drift $\gamma>0$ with residuals
\[
  U_{ij}\leq \Phi_i-\Phi_j-\gamma-r_{ij},
  \qquad r_{\min}:=\min_{i\to j}r_{ij}>0 .
\]
Summing the Lyapunov inequality around a directed cycle telescopes the
potential terms, so every coarse cycle has mean at most
$-\gamma-r_{\min}$. After reserving raw/coarse errors, cutoff variation
$\Delta$, and rational roundoff $M/Q$, the simultaneous strict program surplus
is
\[
  \eta=
  \min\{s_{\min}-\varepsilon_{\rm raw}-\varepsilon_{\rm coarse},
          r_{\min}-\varepsilon_{\rm coarse}\}
  -\Delta-M/Q .
\]
Whenever this number is positive, the finite residual witness supplies the
exact $\eta$ required by the interval-refinement termination theorem. Thus the
remaining positivity target can be checked by finite raw slack residuals and a
finite Lyapunov/Karp drift residual.

The Lyapunov drift residual need not be guessed separately. If the finite
coarse graph has maximum cycle mean $\lambda_{\max}$ and
\[
  \kappa:=-\gamma-\lambda_{\max}>0,
\]
then every directed cycle has mean at most $-\gamma-\kappa$. Applying the
finite negative-cycle-to-potential certificate at the strengthened drift
$\gamma+\kappa$ gives a potential satisfying
\[
  U_{ij}\leq \Phi_i-\Phi_j-\gamma-\kappa .
\]
Thus the Karp margin $\kappa$ is itself a Lyapunov drift residual. Combining
this with the raw slack residual gives
\[
  \eta=
  \min\{s_{\min}-\varepsilon_{\rm raw}-\varepsilon_{\rm coarse},
          \kappa-\varepsilon_{\rm coarse}\}
  -\Delta-M/Q .
\]
This makes the upstream finite target even sharper: prove strict raw slack and
a strict Karp maximum-cycle-mean margin on the genuine Yang--Mills coarse
scale graph.

The same Karp-surplus check is stable under finite interval uncertainty. If a
raw upper bound, coarse raw envelope, or coarse graph edge weight is known with
rigorous interval radius, replace it by its upper endpoint:
\[
  U_{\rm raw}^+=\widehat U_{\rm raw}+\operatorname{err}_{\rm raw},\qquad
  W^+=\widehat W+\operatorname{err}_W,\qquad
  w_{ij}^+=\widehat w_{ij}+\operatorname{err}_{ij}.
\]
If the endpoint data satisfy
\[
  s_{\min}^+:=\min_e\{W^+(e)-U_{\rm raw}^+(e)\}>0,
  \qquad
  \kappa^+:=-\gamma-\lambda_{\max}(w^+)>0,
\]
then monotonicity of edge upper bounds and cycle means shows that every
realization inside the intervals is dominated by the endpoint certificate.
The certified surplus is therefore
\[
  \eta=
  \min\{s_{\min}^+-\varepsilon_{\rm raw}-\varepsilon_{\rm coarse},
          \kappa^+-\varepsilon_{\rm coarse}\}
  -\Delta-M/Q .
\]
So rigorous interval edge evidence with positive worst-case raw slack and
positive worst-case Karp margin already supplies the surplus needed by the
finite validation chain.

Combining this endpoint surplus with a convergence rate gives an explicit
finite validation program. If the interval/Taylor refinement error satisfies
\[
  \operatorname{Err}(h)\leq A h^p
\]
and the endpoint certificate gives surplus $\eta>0$, then every mesh with
\[
  h\leq (\eta/2A)^{1/p}
\]
leaves certified surplus at least $\eta/2$. On a compact $d$-dimensional
parameter box of $\ell^1$ diameter $D$, it suffices to choose
\[
  N\geq D(2A/\eta)^{1/p}
\]
subdivisions per coordinate, giving at most $E N^d$ edge charts. Thus a
worst-case endpoint raw-slack/Karp witness plus a quantitative enclosure
convergence rate yields a concrete finite mesh threshold for validating the
Yang--Mills seed.

This can be packaged as a single endpoint-witness program. One finite object
containing endpoint raw slack, endpoint Karp margin, enclosure convergence
rate, rational denominator, screened bridge norm $C_{\rm SB}$, leakage bridge
constant $K$, packet Lipschitz constant $L_F$, chart anchor count $Q_N$, and
rooting constant $K_A$ determines all deterministic DLCR outputs:
\[
  \ell_x\leq K L_F C_{\rm SB} e^{-m k_N(x)},\qquad
  B_N=K_A Q_N C_{\rm SB},
\]
and hence the effective-action transport constant
\[
  W_1\leq R\sqrt{K_A Q_N C_{\rm SB}}\,a^{\beta/2}.
\]
Together with the mesh bound above, this shows that a single genuine
Yang--Mills endpoint interval witness with common screened/rooting constants
would drive the full finite DLCR certificate program. It is still a conditional
reduction, not a construction of that witness.

The witness can be made into a single finite artifact. A package with schema
\[
  \texttt{yangmills.dlcr.endpoint\_interval\_witness.v1}
\]
and status \texttt{PROVED} must carry the action hash, endpoint raw/coarse
intervals, reserve budgets, refinement model, parameter-domain data, scalar
complexity, and screened constants. If the action hash matches the expected
Yang--Mills action, the verifier extracts all fields and applies the endpoint
interval witness program above. Thus the final deterministic target is one
hash-matched endpoint interval witness artifact; producing that artifact would
supply the finite mesh, chart bound, local leakage constant, effective-action
constant, and transport constant.

The raw Taylor-atlas target and the endpoint interval target are in fact one
finite pipeline. Suppose a package with schema
\[
  \texttt{yangmills.dlcr.raw\_edge\_taylor\_atlas.v1}
\]
has status \texttt{PROVED} and its action hash matches the ambient
Yang--Mills action. The raw-edge schema certificate extracts certified raw
upper bounds $U_e$, canonical coarse envelopes $W_e$, canonical coarse
transition weights $w_f$, and the parameter box. Given nonnegative interval
radii $r_{\rm raw},r_W,r_w$, reserve budgets, a refinement model
$\operatorname{Err}(h)\leq A h^p$, denominator $Q$, scalar complexity $M$,
and common screened constants, define an endpoint witness by using the
interval centers $U_e,W_e,w_f$ and radii $r_{\rm raw},r_W,r_w$. The endpoint
verifier then checks the worst-case upper endpoints
\[
  U_e^+=U_e+r_{\rm raw},\qquad W_e^+=W_e+r_W,\qquad
  w_f^+=w_f+r_w
\]
and computes
\[
  \eta=\min\{s_{\min}^+-\varepsilon_{\rm raw}-\varepsilon_{\rm coarse},
             \kappa^+-\varepsilon_{\rm coarse}\}-\Delta-\frac{M}{Q}.
\]
If this number is positive, the endpoint interval witness schema theorem
applies and produces the finite DLCR program constants. Thus a proved raw
Taylor-atlas/Karp package plus explicit radii and screened constants
automatically constructs the endpoint interval witness; no independent
endpoint artifact remains to be invented. If the radii or reserves destroy the
raw slack or Karp margin, the verifier rejects the conversion.

The local $C^2$ enclosure target can itself be generated by a finite
sectorized polymer majorant package. On a sector chart, suppose the raw edge
log norm is controlled by a value center $v_0$ with radius $r_v$, gradient
centers $g_i$ with radii $r_i$, and a nonnegative Hessian $\ell^1$ majorant
$H$. The corresponding local enclosure is simply
\[
  [v_0-r_v,v_0+r_v],\qquad [g_i-r_i,g_i+r_i],\qquad [0,H].
\]
Thus the next analytic gate can be formulated without hidden functional
analysis: prove a finite sector decomposition of the Yang--Mills polymer
activities and rigorous value/gradient/Hessian majorants on each sector. The
downstream verifier still checks support, cover, domination, and Karp
negativity, so sector majorants do not assume the screened norm or the mass
gap.

The jet-atlas target itself can be fed by ordinary local $C^2$ enclosure data.
Suppose a proved local enclosure package gives, on each raw edge chart, a
base-point value interval $[v^-,v^+]$, gradient intervals
$[g_i^-,g_i^+]$, and a nonnegative Hessian $\ell^1$ interval $[H^-,H^+]$.
The canonical jet data are
\[
  v_0={v^-+v^+\over2},\quad r_v={v^+-v^-\over2},\qquad
  H_0={H^-+H^+\over2},\quad r_H={H^+-H^-\over2},
\]
with the same gradient intervals, base point, and parameter box. Thus the
Ward-projected jet atlas is no longer a mysterious object: it is exactly the
finite packaging of rigorous local value, gradient, and Hessian interval
estimates, plus raw/coarse support and Karp data. The difficult Yang--Mills
work is to prove those local $C^2$ enclosures and keep their margins uniform
through the cutoff and volume limits.

There is now one further finite synthesis layer before the raw-edge package.
A \emph{Ward-projected jet atlas} records, on each raw edge chart, a parameter
box, a base point, an interval value at the base point, interval gradients,
and an interval Hessian bound. From such data one constructs an ordinary
Taylor-atlas chart by replacing the value by its upper endpoint, using the
gradient interval midpoint as the affine coefficient, and adding the explicit
gradient uncertainty budget
\[
  E_{\nabla}=\sum_i {g_i^+-g_i^-\over 2}\,R_i
\]
to the base value, where $R_i$ is the coordinate radius of the chart box. For
every admissible gradient $g_i=g_i^0+e_i$ and every point in the box,
\[
  \sum_i e_i(x_i-x_{0,i})\leq \sum_i |e_i|R_i\leq E_{\nabla}.
\]
The Hessian bound is replaced by its upper endpoint, so Taylor's theorem gives
a valid outward edge envelope. The raw-edge schema verifier is then run on the
synthesized package and still checks support, cover, domination, and Karp
negativity. Hence a proved Ward-projected jet atlas with matching action hash
is enough to produce the raw Taylor-atlas/Karp evidence; the remaining
Yang--Mills analytic gate becomes the construction of those interval jets and
coarse Karp data, rather than a separate endpoint-witness problem.

A useful refinement certificate lets one verify estimates on a finer atlas and
then collapse them to coarser analytic regimes. Let fine charts $B_c$ refine
coarse charts $B_C$ through an assignment $\rho(c)=C$, with
$B_c\subseteq B_C$, and suppose the Taylor-atlas calculation has certified
\[
  w_e(x)\leq U_c,\qquad x\in B_c .
\]
Define the coarse envelope
\[
  V_C:=\max_{\rho(c)=C} U_c .
\]
If $x$ lies in the portion of $B_C$ covered by the assigned fine charts, then
$x\in B_c$ for some $c$ with $\rho(c)=C$, hence
\[
  w_e(x)\leq U_c\leq V_C .
\]
Thus the fine atlas may be replaced by coarse envelope values $V_C$, and
$\max_C V_C=\max_c U_c$ preserves the global edge bound used by the
coarse-graph certificate. This makes adaptive subdivision compatible with a
small finite state graph, but it does not prove the fine Taylor estimates or
the refinement relation for the actual Yang--Mills regimes.

The finite data needed from Yang--Mills can therefore be packaged as a raw-edge
evidence schema. Such a package must give the parameter domain, the finite raw
edges, Taylor charts with base values, gradients, and Hessian bounds, optional
grid cover data, a raw-to-coarse state map, coarse edge envelopes, and the
constants used in the Karp cycle-mean test. If the package has status
\texttt{PROVED}, all chart boxes lie inside the parameter domain, declared grid
cells match the Taylor charts, the raw-to-coarse state partition is supported
exactly on the raw edge endpoints, every endpoint image is a vertex of the
canonical coarse graph, every raw edge maps to a coarse edge, every
Taylor-atlas edge envelope is dominated by its coarse edge, and the coarse
Karp margin is positive, then the existing Taylor-atlas coarse graph
certificate gives the screened bridge norm. This is a schema soundness theorem:
it identifies the exact finite proof artifact that would discharge the edge
gate, while the construction of that artifact from the Yang--Mills measure is
still open.

The finite witness also canonically extracts explicit edge-error budgets. Let
\[
  s_{\min}=\min_e\{W_{\rm coarse}(\pi(e))-U_e\}
\]
be the minimum slack between the certified raw upper bound $U_e$ and its
canonical coarse envelope, and let
\[
  m_{\rm cyc}=-\gamma-\lambda_{\max}
\]
be the coarse Karp cycle margin. If later analytic estimates increase every
raw edge bound by at most $\varepsilon_{\rm raw}$ and every coarse edge bound
by at most $\varepsilon_{\rm coarse}$, then the same certificate survives
provided
\[
  \varepsilon_{\rm raw}+\varepsilon_{\rm coarse}\leq s_{\min},
  \qquad
  \varepsilon_{\rm coarse}\leq m_{\rm cyc}.
\]
The first inequality preserves every raw-to-coarse envelope, and the second
uses monotonicity of cycle means under edge perturbations. This turns future
Yang--Mills edge estimates into a margin-fitting problem with explicit error
budgets. In particular, any passing raw-edge evidence package yields the
finite acceptance region
\[
  \{(\varepsilon_{\rm raw},\varepsilon_{\rm coarse}):
  \varepsilon_{\rm raw}+\varepsilon_{\rm coarse}\leq s_{\min},\quad
  \varepsilon_{\rm coarse}\leq m_{\rm cyc}\}.
\]
The verifier reports this region, together with a balanced admissible point
obtained by taking
\[
  \varepsilon_{\rm coarse}=\min\{s_{\min}/2,m_{\rm cyc}\},\qquad
  \varepsilon_{\rm raw}=s_{\min}-\varepsilon_{\rm coarse}.
\]
Thus the remaining analytic construction need not hit exact edge values: it
only has to produce raw and coarse error bars fitting inside this certified
finite region.

The same raw-edge package is stable under refinement. If a refined package
shares the same schema, parameter domain, and state partition, and if its
screened bridge constant does not increase while its coarse Karp margin does
not decrease, then it is a safe replacement for the coarse package. That
monotonicity is useful because adaptive subdivision should never make a valid
finite witness worse, but it still does not construct the Yang--Mills witness.

Once the finite witness is also tied to the ambient action by the provenance
hash, the deterministic output layer is completely compositional. A passing
raw-edge package gives a screened bridge constant $C_{\rm SB}$. The
screened-norm-to-leakage certificate gives
\[
  \ell_x(a,L)\leq K L_F C_{\rm SB} e^{-m k_N(x)}
\]
for any supplied curvature-to-leakage constant $K$, packet Lipschitz bound
$L_F$, and screening rate $m>0$. The screened-norm-to-effective-action
certificate gives, for $Q_N$ chart anchors and bridge-to-action constant
$K_A$,
\[
  \|S_a-S\|_\infty\leq K_A Q_N C_{\rm SB}a^\beta,
  \qquad
  W_1\leq R\sqrt{K_A Q_N C_{\rm SB}}\,a^{\beta/2}.
\]
Thus a single hash-matched raw Taylor-atlas/Karp evidence package would
simultaneously supply the screened norm, local exponential leakage constant,
and finite-chart effective-action transport constants. This is the most compact
finite certificate target left in the route. The analytic burden is not hidden:
one must still construct the genuine Yang--Mills raw-edge witness and prove the
curvature bridge/rooting normalizations and cutoff-uniformity assumptions.

\subsection{No current leakage}
Finally, the finite charts must be compatible with the countable exhaustion.
The fixed-packet cylindrical part is discharged. If $\Gamma_N$ contains all
edges used by a packet observable $O$, then $O=O_N\circ\pi_N$ for the projection
onto the $\Gamma_N$ edge variables. Changing all complement edges changes $O$ by
exactly $0$.

The remaining leakage problem is uniform and analytic. One still needs an
estimate compatible with the continuum and countable-exhaustion limits:
The proposed leakage estimate is
\[
  \lim_{N\to\infty}\limsup_{a\to0,L\to\infty}\epsilon_N(a,L)=0,
\]
where $\epsilon_N$ bounds the contribution of current outside the packet chart.
This is the missing uniformity needed to connect the finite current comparison
to the existing countable localization layer.

Equivalently, it is enough to construct a deterministic tail schedule
$\rho_N\to0$ such that
\[
  \limsup_{a\to0,L\to\infty}\epsilon_N(a,L)\leq\rho_N .
\]
Then
\[
  0\leq
  \limsup_{a\to0,L\to\infty}\epsilon_N(a,L)
  \leq\rho_N
\]
for each $N$, and sending $N\to\infty$ gives the required no-leakage limit. The
new finite content is only this tail-schedule implication; the analytic
Yang--Mills localization estimate producing such a $\rho_N$ remains open.

A standard exponential localization estimate would produce the required
schedule explicitly. If the packet exhaustion has radii $R_N\to\infty$ and
\[
  \limsup_{a\to0,L\to\infty}\epsilon_N(a,L)
  \leq
  C e^{-mR_N},
  \qquad C\geq0,\quad m>0,
\]
then one may take
\[
  \rho_N=C e^{-mR_N}.
\]
Since $R_N\to\infty$, $\rho_N\to0$, and the previous criterion applies. This
again only converts a localization theorem into the DLCR leakage artifact; it
does not prove the Yang--Mills localization estimate.

There is also a summable-shell route. Suppose the complement of the $N$-th
packet chart is decomposed into shells with nonnegative leakage bounds $s_k$
and
\[
  \sum_{k\geq1}s_k<\infty .
\]
If the continuum and infinite-volume limsup leakage satisfies
\[
  \limsup_{a\to0,L\to\infty}\epsilon_N(a,L)
  \leq
  \sum_{k\geq N}s_k ,
\]
then one may take
\[
  \rho_N=\sum_{k\geq N}s_k .
\]
The tails of a convergent nonnegative series tend to zero, so this again gives
a deterministic no-leakage schedule. As with the exponential route, this is
only a certificate conversion: the Yang--Mills input would be the actual
uniform summable shell estimate.

A common way to produce such shell bounds is polynomial shell growth combined
with exponential local decay. If the $k$-th complement shell has at most
$V(k+1)^p$ atoms, with $V\geq0$ and $p\geq0$, and each atom contributes at most
$C e^{-mk}$ to the continuum/infinite-volume limsup leakage, with $C\geq0$ and
$m>0$, then
\[
  s_k=VC(k+1)^p e^{-mk}
\]
is summable. Indeed, for $p>0$ the elementary bound
\[
  x^p\leq \left(\frac{p}{e\alpha}\right)^p e^{\alpha x},
  \qquad \alpha=m/2,
\]
gives
\[
  (k+1)^p e^{-mk}
  \leq
  A_{p,m}e^{-mk/2},
  \qquad
  A_{p,m}=e^{m/2}\left(\frac{2p}{em}\right)^p ,
\]
while for $p=0$ one takes $A_{0,m}=1$. Therefore
\[
  \sum_{k\geq N}s_k
  \leq
  \frac{VCA_{p,m}e^{-mN/2}}{1-e^{-m/2}} .
\]
This proves summability with explicit constants. The remaining Yang--Mills
work would be to prove the actual polynomial shell-volume estimate and the
uniform exponential local leakage estimate.

For finite lattice packet supports, the polynomial shell-volume part is itself
elementary. If a packet support has $Q$ anchors in $\mathbb Z^d$, the
$\ell^\infty$ shell of radius $k$ around one anchor has exact cardinality
\[
  1\quad(k=0),\qquad
  (2k+1)^d-(2k-1)^d\quad(k\geq1).
\]
For $k\geq1$, the mean value theorem gives
\[
  (2k+1)^d-(2k-1)^d
  \leq
  2d(2k+1)^{d-1}
  \leq
  2d\,3^{d-1}(k+1)^{d-1}.
\]
The same bound trivially covers $k=0$, and a union bound over the $Q$ anchors
gives
\[
  |\mathrm{shell}_k|
  \leq
  Q\,2d\,3^{d-1}(k+1)^{d-1}.
\]
Thus the leakage side has been reduced further: for the lattice shells one may
take $V=Q\,2d\,3^{d-1}$ and $p=d-1$. What remains is the genuine Yang--Mills
estimate that shell leakage decays like $Ce^{-mk}$ with constants uniform in
the continuum and infinite-volume limits.

The exact local estimate needed can be stated atomwise. Let $S_N$ be the packet
support and let
\[
  k_N(x)=\operatorname{dist}_\infty(x,S_N)
\]
be the shell index of a lattice atom $x$. If $\ell_x(a,L)\geq0$ bounds the
change in the packet chart caused by the outside atom $x$, then the leakage is
controlled by
\[
  \epsilon_N(a,L)
  \leq
  \sum_{k\geq N}\sum_{x:\,k_N(x)=k}\ell_x(a,L).
\]
Thus a uniform local exponential estimate
\[
  \limsup_{a\to0,L\to\infty}\ell_x(a,L)
  \leq
  C e^{-m k_N(x)},\qquad C<\infty,\quad m>0,
\]
would give
\[
  \limsup_{a\to0,L\to\infty}\epsilon_N(a,L)
  \leq
  \frac{Q\,2d\,3^{d-1}C A_{d-1,m}e^{-mN/2}}{1-e^{-m/2}},
\]
which is the required deterministic leakage tail schedule. The repository now
records this as a conditional leakage evidence artifact: all finite counting
and summation constants are present, but the artifact is not promoted because
the Yang--Mills local exponential leakage estimate itself is still missing.
The verifier now requires this completely established estimate as its own action-hash-matched
artifact, with schema
\[
  \texttt{yangmills.dlcr.local\_exponential\_leakage.v1}.
\]
Thus even a promoted tail-schedule file cannot discharge no-current-leakage
unless the local estimate
\[
  \limsup_{a\to0,L\to\infty}\ell_x(a,L)
  \leq C e^{-m k_N(x)}
\]
is separately proved with constants $C<\infty$ and $m>0$ and without appealing
to the mass gap being proved.

\subsection{True analytic inputs}
The preceding reductions leave three non-formal Yang--Mills estimates. They are
not supplied by the existing broad tightness, moment-bound, operator-convergence,
or semigroup evidence. To promote this DLCR route, the repository now requires
three explicit finite-chart artifacts, all pinned to the same Wilson-action
specification:
\begin{enumerate}
\item a state chart transport estimate
\[
  \mathbf E\,d_\infty(X_a,Y)\leq D_Na^{\eta_N}
\]
for every fixed Wilson-loop packet $W_N$;
\item translated chart transport estimates
\[
  \mathbf E\,d_\infty(X_{a,q},Y_q)\leq D_{N,q}a^{\eta_{N,q}}
\]
for every rational-time entry in $F_N$;
\item a leakage tail schedule
\[
  \limsup_{a\to0,L\to\infty}\epsilon_N(a,L)\leq\rho_N,\qquad \rho_N\to0.
\]
\end{enumerate}
The two transport artifacts may be proved either by constructing couplings
directly or by a finite Kantorovich-dual certificate. Namely, on a compact
finite chart it suffices to give an $\varepsilon(a)$-net
$\{\psi_i\}$ for the anchored $1$-Lipschitz test class and prove
\[
  \left|\mathbf E_{\mu_a}\psi_i-\mathbf E_{\mu_{\mathrm{OS}}}\psi_i\right|
  \leq A a^\eta
\]
for every net test. Kantorovich--Rubinstein duality then gives
\[
  W_1(\mu_a,\mu_{\mathrm{OS}})
  \leq A a^\eta+2\varepsilon(a).
\]
If $\varepsilon(a)\leq Ea^\gamma$, this is a transport bound with constant
$A+2E$ and exponent $\min(\eta,\gamma)$.
Another accepted route is relative entropy. If a finite chart has diameter
$R$ and
\[
  \operatorname{KL}(\mu_a\|\mu_{\mathrm{OS}})\leq K a^\kappa,
\]
then Pinsker's inequality and the maximal coupling bound give
\[
  W_1(\mu_a,\mu_{\mathrm{OS}})
  \leq R\sqrt{K/2}\,a^{\kappa/2}.
\]
For $d_\infty$ on $SU(3)^m$ representatives one may take $R\leq2$, hence the
constant is $\sqrt{2K}$.
This entropy estimate follows in turn from uniform convergence of the finite
chart effective actions. If
\[
  d\mu_a=Z_a^{-1}e^{-S_a}\,d\lambda,\qquad
  d\mu_{\mathrm{OS}}=Z^{-1}e^{-S}\,d\lambda
\]
on a common finite reference measure and
$\|S_a-S\|_\infty\leq B a^\beta$, then
$e^{-Ba^\beta}Z\leq Z_a\leq e^{Ba^\beta}Z$. Hence the log density ratio is
bounded above by $2Ba^\beta$, so
\[
  \operatorname{KL}(\mu_a\|\mu_{\mathrm{OS}})\leq 2Ba^\beta
\]
and therefore
\[
  W_1(\mu_a,\mu_{\mathrm{OS}})
  \leq R\sqrt{B}\,a^{\beta/2}.
\]
For $d_\infty$ on $SU(3)^m$, this gives constant $2\sqrt B$.
Finally, the uniform effective-action estimate follows from a summable
finite-chart polymer or cumulant remainder. If
\[
  S_a-S=\sum_\Gamma c_\Gamma(a)K_\Gamma,\qquad
  \|K_\Gamma\|_\infty\leq1,
\]
and
\[
  |c_\Gamma(a)|\leq A_\Gamma a^{\beta_\Gamma},\qquad
  \beta_\Gamma\geq\beta>0,\qquad
  \sum_\Gamma A_\Gamma\leq B,
\]
then
\[
  \|S_a-S\|_\infty\leq Ba^\beta.
\]
Thus a packet-scale polymer/cumulant summability theorem with explicit
$B,\beta$ would supply the effective-action estimate, hence the entropy and
transport estimates above.

The finite interaction-degree step is discharged for finite-range lattice
polymer atoms. If atoms are indexed by a subset of $\mathbb Z^d$ and two atoms
interact only when their $\ell^\infty$ distance is at most $r$, then the
interaction graph has degree
\[
  \Delta\leq (2\lceil r\rceil+1)^d-1 .
\]
Indeed, every neighbor of a fixed atom lies in the cube of side length
$2\lceil r\rceil+1$ centered at that atom, excluding the atom itself.

The finite shape-counting step is discharged for any such bounded-degree
polymer interaction graph. If the chart support has $Q_N$ anchors, then
connected polymers of size $n$ touching the chart support are bounded in number
by
\[
  Q_N\mu^n,
  \qquad
  \mu=\max(1,2e\Delta).
\]
Indeed, for one fixed root, a connected size-$n$ polymer is overcounted by a
rooted subtree of the universal covering tree. A preorder encoding chooses
$n-1$ occupied child slots among at most $\Delta n$ slots, so
\[
  \binom{\Delta n}{n-1}
  \leq
  \left(\frac{e\Delta n}{n-1}\right)^{n-1}
  \leq
  (2e\Delta)^{n-1}
\]
for $n\geq2$, while $n=1$ is trivial. The union bound over the $Q_N$ anchors
gives the displayed estimate.

Suppose, in addition, that the corresponding activities obey
\[
  |z_\Gamma(a)|\leq A_0a^\beta q^n,\qquad \mu q<1 .
\]
Equivalently, an exponential activity estimate
\[
  |z_\Gamma(a)|\leq A_0a^\beta e^{-\sigma n}
\]
is sufficient whenever $\sigma>\log\mu$, by taking $q=e^{-\sigma}$.
Likewise, if a verified cluster-decay certificate supplies a ratio $r$ with
$0<r<1/\mu$, then $\sigma=-\log r$ and the same threshold follows.
Then grouping by $n$ gives the geometric estimate
\[
  \sum_\Gamma |z_\Gamma(a)|
  \leq
  Q_NA_0a^\beta\sum_{n\geq1}(\mu q)^n
  =
  Q_NA_0\frac{\mu q}{1-\mu q}\,a^\beta .
\]
Thus the polymer-remainder certificate applies with
$B=Q_NA_0\mu q/(1-\mu q)$. This proves the bounded-degree
shape-growth estimate and the finite Kotecky--Preiss-style
activity-to-summability implication. The actual Yang--Mills polymer expansion,
the finite-range polymer atom identification, and the excess-decay estimate
$\sigma>\log\mu$ remain analytic inputs.

Combining the finite-range degree, shape-growth, cluster-decay, polymer
remainder, entropy, and transport certificates gives a single usable transport
package. If a finite-chart Yang--Mills expansion supplies finite-range atoms in
$\mathbb Z^d$, interaction radius $r$, chart support size $Q_N$, activity
prefactor $A_0a^\beta$, and a verified cluster ratio $\rho<1/\mu$, with
\[
  \mu=\max\!\left(1,2e\big((2\lceil r\rceil+1)^d-1\big)\right),
\]
then
\[
  B=Q_NA_0\frac{\mu\rho}{1-\mu\rho}
\]
gives
\[
  \operatorname{KL}(\mu_a\|\mu_{\mathrm{OS}})
  \leq 2B a^\beta,
  \qquad
  W_1(\mu_a,\mu_{\mathrm{OS}})
  \leq R\sqrt B\,a^{\beta/2}.
\]
This packages the finite constants only; the Yang--Mills expansion and the
ratio $\rho<1/\mu$ remain the genuine analytic target.

The verifier now also records a single unified route for this target. A future
artifact with schema
\[
  \texttt{yangmills.dlcr.finite\_chart\_polymer\_cluster\_expansion.v1}
\]
would be sufficient to populate the three remaining DLCR evidence files only
if it proves all of the following output estimates with constants and the same
action hash:
\[
  \mathbf E\,d_\infty(X_a,Y)\leq D_Na^{\eta_N},
\]
\[
  \mathbf E\,d_\infty(X_{a,q},Y_q)\leq D_{N,q}a^{\eta_{N,q}},
\]
and
\[
  \ell_x(a,L)\leq C e^{-m k_N(x)}.
\]
The first output feeds state chart transport, the second feeds translated chart
transport, and the third feeds the leakage-tail schedule through the shell
estimates above. This unified route is only a promotion criterion: the current
artifact is conditional and is not accepted as a substitute for the three
proved output artifacts.

One proposed way to prove such a unified artifact is the Ward-screened cavity
route. For an outside atom $x$, introduce a cavity interpolation and decompose
\[
  \partial_t S_{t,x}
  =
  \nabla_{\mathrm{gauge}}G_{t,x}+B^{\mathrm{curv}}_{t,x}.
\]
For a finite compact chart law
\[
  d\mu_t=Z_t^{-1}e^{-S_t}d\lambda,
\]
the finite cavity interpolation identity gives
\[
  \frac{d}{dt}\mathbf E_t[F]
  =
  -\operatorname{Cov}_t(F,\dot S_t).
\]
If $\dot S_t=XG_t+B^{\mathrm{curv}}_t$ and the finite Ward identity gives
$\operatorname{Cov}_t(F,XG_t)=0$, then
\[
  \frac{d}{dt}\mathbf E_t[F]
  =
  -\operatorname{Cov}_t(F,B^{\mathrm{curv}}_t).
\]
The required Ward cancellation lemma is
\[
  \operatorname{Cov}_{t,x}
  \bigl(F(J_N),\nabla_{\mathrm{gauge}}G_{t,x}\bigr)=0
\]
for gauge-invariant packet tests $F$. This cancellation is a finite compact
chart identity: Haar integration by parts gives
$\int X(H)\,d\mu=0$ for gauge generators $X$ and invariant finite-chart
densities, and $X(F)=0$ for gauge-invariant packet tests, hence
$\operatorname{Cov}(F,XG)=0$. What remains is to prove that the cavity
derivative actually decomposes into such gauge-generator terms plus a curvature
bridge term. The remaining curvature bridge term must then admit connected
cluster weights $w_\Gamma$ satisfying the screened norm bound
\[
  \sup_{N,x}
  e^{m k_N(x)}
  \sum_{\Gamma:\,x\leftrightarrow S_N}
  |w_\Gamma|e^{\sigma|\Gamma|}
  <\infty,
  \qquad \sigma>\log\mu,\quad m>0.
\]
A one-shell contraction
\[
  \|B_{r+1}\|_{\mathrm{SB}}
  \leq
  \theta\|B_r\|_{\mathrm{SB}}+Ca^\beta e^{-\sigma r},
  \qquad \theta<1,
\]
together with uniform continuum and volume control, would imply the local
exponential leakage estimate and the finite-chart effective-action remainders.
The deterministic part of this implication is now separated. If nonnegative
shell bridge norms $b_r$ satisfy
\[
  b_{r+1}\leq\theta b_r+Ca^\beta e^{-\sigma r},
  \qquad 0<a\leq1,\quad 0\leq\theta<1,
\]
then for every screening rate $m$ with
\[
  \theta e^m<1,\qquad m<\sigma,
\]
one has
\[
  \sup_{r\geq0}e^{mr}b_r
  \leq
  b_0+
  \frac{Ce^m}{(1-\theta e^m)(1-e^{m-\sigma})}.
\]
Indeed, iteration gives
\[
  b_r\leq
  \theta^r b_0+
  Ca^\beta\sum_{j=0}^{r-1}\theta^{r-1-j}e^{-\sigma j},
\]
and multiplication by $e^{mr}$ bounds the forced term by the product of the
two geometric series
$\sum_{\ell\geq0}(\theta e^m)^\ell$ and
$\sum_{j\geq0}(e^{m-\sigma})^j$. Thus a genuine Yang--Mills proof of the
one-shell recurrence, with uniform $b_0,\theta,C,\sigma,\beta$, would give the
screened bridge norm with the displayed constant.

A second route avoids reducing the bridge propagation to nearest-shell
recurrences. Let $b_x$ be the Ward-projected bridge norm rooted at an outside
atom $x$ and suppose the curvature bridge analysis gives an influence
inequality
\[
  b_x\leq f_x+\sum_y A_{xy}b_y,\qquad A_{xy}\geq0,
\]
with forcing
\[
  f_x\leq C a^\beta e^{-\sigma k_N(x)} .
\]
For a screening rate $0<m<\sigma$, define the weighted influence norm
\[
  \theta_m=
  \sup_x\sum_y A_{xy}\exp\{m(k_N(x)-k_N(y))\}.
\]
If $\theta_m<1$, then
\[
  \sup_x e^{m k_N(x)}b_x
  \leq
  \frac{C}{1-\theta_m}a^\beta .
\]
Indeed, multiply the influence inequality by $e^{mk_N(x)}$, take the supremum
over $x$, and absorb the term $\theta_m\sup_y e^{mk_N(y)}b_y$ to the left.
This Ward-projected Dobrushin resolvent certificate is deterministic. It gives
a second possible path to the screened bridge norm, with
\[
  C_{\mathrm{SB}}=\frac{C}{1-\theta_m},
\]
provided the genuine Yang--Mills bridge analysis proves the weighted influence
bound $\theta_m<1$ uniformly in the continuum and infinite-volume limits.
There is also an iterated version. Let $T_m$ be the weighted influence operator
with kernel $A_{xy}\exp\{m(k_N(x)-k_N(y))\}$. If
\[
  \theta_{m,p}:=\|T_m^p\|_{\infty}<1
\]
for some integer $p\geq1$, then the same positive-resolvent argument gives
\[
  \sup_x e^{m k_N(x)}b_x
  \leq
  C a^\beta
  \frac{1+\theta_m+\cdots+\theta_m^{p-1}}{1-\theta_{m,p}},
  \qquad \theta_m=\|T_m\|_\infty .
\]
This can be useful if the one-step weighted norm is not contractive but a
finite Ward-projected iterate is. It is still only a deterministic certificate:
the Yang--Mills task is to prove such an iterated weighted influence
contraction with constants independent of the cutoff limits.
The atom-level influence check can also be compressed to finitely many blocks.
Partition the outside atoms into blocks $I_1,\ldots,I_q$ and suppose a
nonnegative matrix $B$ satisfies
\[
  \sup_{x\in I_i}\sum_{y\in I_j}
  A_{xy}\exp\{m(k_N(x)-k_N(y))\}\leq B_{ij}.
\]
Then the weighted atom influence operator $T_m$ satisfies
\[
  \|T_m^p\|_\infty\leq \|B^p\|_\infty
\]
for every $p\geq1$. Indeed, the block sup vector of $T_mu$ is bounded
componentwise by $B$ times the block sup vector of $u$, and iteration gives the
claim. Consequently a contractive finite block iterate $\|B^p\|_\infty<1$
implies the atom-level iterated Dobrushin certificate. This converts the
analytic target from an atom-by-atom estimate into finitely many block
influence bounds, still with the same uniformity requirements.

The final deterministic conversion from screened norm to local leakage is also
separated. Suppose the curvature bridge expansion bounds the leakage of an
outside atom $x$ by
\[
  \ell_x(a,L)
  \leq
  K L_F
  \sum_{\Gamma:\,x\leftrightarrow S_N}|w_\Gamma|,
\]
where $K$ is the bridge-to-leakage normalization and $L_F$ is the packet-test
Lipschitz constant. If
\[
  \sup_x e^{m k_N(x)}
  \sum_{\Gamma:\,x\leftrightarrow S_N}
  |w_\Gamma|e^{\sigma|\Gamma|}
  \leq C_{\mathrm{SB}},
  \qquad m>0,\quad \sigma\geq0,
\]
then $e^{\sigma|\Gamma|}\geq1$ gives
\[
  \ell_x(a,L)
  \leq
  K L_F C_{\mathrm{SB}} e^{-m k_N(x)}.
\]
Thus a proved screened bridge norm would produce the exact local exponential
leakage artifact required above, with local leakage constant
$K L_F C_{\mathrm{SB}}$ and rate $m$.

The same screened norm also gives the finite-chart effective-action constants
needed for the transport side. Suppose every chart-touching remainder cluster
can be assigned to one anchor in the finite chart support $S_N$, with
$|S_N|=Q_N$, and its contribution to the effective-action difference is bounded
by $K_Aa^\beta |w_\Gamma|$. Since $e^{\sigma|\Gamma|}\geq1$ and the shell
distance from a chart anchor to $S_N$ is zero, the rooted cluster sum at each
anchor is at most $C_{\mathrm{SB}}$. A union bound over the $Q_N$ anchors gives
\[
  \|S_a-S\|_\infty
  \leq
  K_A Q_N C_{\mathrm{SB}} a^\beta .
\]
The finite effective-action certificate therefore gives
\[
  \operatorname{KL}(\mu_a\|\mu_{\mathrm{OS}})
  \leq
  2K_A Q_N C_{\mathrm{SB}} a^\beta
\]
and, on a chart of diameter $R$,
\[
  W_1(\mu_a,\mu_{\mathrm{OS}})
  \leq
  R\sqrt{K_A Q_N C_{\mathrm{SB}}}\,a^{\beta/2}.
\]
The translated chart has the same formula with $Q_N,R$ replaced by the
translated packet constants.

The verifier records this Ward-screened cavity route as a conditional framework
with six sublemmas. The finite cavity interpolation and finite Ward
cancellation sublemmas are proved, and the formal one-shell-recurrence-to-
screened-norm, Dobrushin/block-compression, multiscale transport-flow,
log-budget, sparse bad-scale, sliding-window, finite-profile, and
negative-average/Lyapunov/finite-state/cycle-mean budgets,
finite-state graph duality and maximum-cycle-mean checks,
finite scale-graph-to-screened-norm packaging and interval robustness,
screened-norm-to-local-leakage, and
screened-norm-to-effective-action reductions are proved; the cavity
decomposition, bridge cluster expansion, the actual Yang--Mills contraction or
negative log-budget/sparse bad-scale/window-density/finite-profile/negative-
average/Lyapunov/finite-state/cycle-mean estimate, and uniform
continuum/volume control remain conditional.

Until these estimates are proved with constants, the exact shadow coupling,
reflection-time intertwiner, and no-current-leakage lemmas remain conditional.

\subsection{Sector cluster derivative majorants}
The latest finite reduction moves the Ward-screened jet-atlas target one layer
closer to a standard polymer expansion theorem. A proved artifact of schema
\[
  \texttt{yangmills.dlcr.sector\_cluster\_derivative\_majorants.v1}
\]
with the repository action hash, finite sector charts, signed truncated cluster
value and first-derivative jets, nonnegative Hessian $\ell^1$ contributions,
absolute value/gradient/Hessian tail budgets, positive KP slack, raw/coarse
support data, and coarse Karp constants canonically constructs the
\texttt{yangmills.dlcr.polymer\_sector\_c2\_majorants.v1} package.
On a sector $C$ the construction is
\[
 v_C=\sum_{\Gamma\in T_C}v_\Gamma,\qquad
 r_C=\sum_{\Gamma\notin T_C}|v_\Gamma|,
\]
\[
 g_{C,i}=\sum_{\Gamma\in T_C}\partial_i v_\Gamma,\qquad
 \rho_{C,i}=\sum_{\Gamma\notin T_C}|\partial_i v_\Gamma|,
\]
and
\[
 H_C=\sum_{\Gamma\in T_C}\|D^2v_\Gamma\|_1+
     \sum_{\Gamma\notin T_C}\|D^2v_\Gamma\|_1 .
\]
The triangle inequality gives the value and gradient intervals
$[v_C-r_C,v_C+r_C]$ and $[g_{C,i}-\rho_{C,i},g_{C,i}+\rho_{C,i}]$; Hessian
subadditivity gives the $\ell^1$ majorant $H_C$. The verifier checks coordinate
agreement, nonnegative tails, positive KP slack, cluster-count domination,
action-hash agreement, and then applies the already proved
polymer-sector-to-local-$C^2$ and local-$C^2$-to-jet-atlas converters. This is a
finite reduction: the genuine Yang--Mills sector cluster expansion and its
derivative-tail bounds remain open.

\subsection{Geometric derivative-tail envelope}
A further finite reduction replaces raw derivative-tail declarations by a
geometric order envelope. A proved artifact of schema
\[
  \texttt{yangmills.dlcr.geometric\_cluster\_tail\_envelope.v1}
\]
contains truncated sector cluster jets through order $N_C$ and constants
$A_v,r_v,A_i,r_i,A_H,r_H$ satisfying $0\le r_\bullet<1$. The omitted tails are
then computed rather than assumed:
\[
 \sum_{n>N_C} A_v r_v^n={A_v r_v^{N_C+1}\over 1-r_v},\qquad
 \sum_{n>N_C} A_i r_i^n={A_i r_i^{N_C+1}\over 1-r_i},
\]
and
\[
 \sum_{n>N_C} A_H r_H^n={A_H r_H^{N_C+1}\over 1-r_H}.
\]
If the sector shape-growth constant $\mu_C$ and activity ratio $\rho_C$ obey
$\mu_C\rho_C<1$, the KP slack is $1-\mu_C\rho_C>0$. The verifier constructs the
sector-cluster derivative-majorant package with these computed tails and then
applies the sector-cluster-to-polymer-sector converter. Thus the next analytic
target is a geometric derivative-tail theorem for the actual Ward-screened
Yang--Mills sector expansion, with enough downstream raw/Karp margin.

\subsection{Finite shell-ratio certificate}
The geometric tail envelope itself is now reduced to a finite shell-ratio
certificate. On each sector, suppose the value, each first derivative, and the
Hessian $\ell^1$ norm have consecutive shell majorants $B_n$ in the tail regime
and a ratio $0\le r<1$ such that
\[
  B_{n+1}\le rB_n .
\]
Anchoring at shell $n_0$ gives the geometric envelope
\[
  B_n\le B_{n_0}r^{n-n_0}
      =\left({B_{n_0}\over r^{n_0}}\right)r^n
      \qquad(n\ge n_0),
\]
with the evident zero-ratio convention. If one activity ratio $\rho$ dominates
the value, gradient, and Hessian ratios and satisfies $\mu\rho<1$ for the
sector shape-growth constant $\mu$, then the KP slack is $1-\mu\rho>0$. The
verifier turns this finite ratio certificate into the geometric
cluster-tail-envelope artifact and then applies the geometric-tail and
sector-cluster converters. The remaining analytic issue is therefore a finite
Yang--Mills shell-ratio theorem for the Ward-screened sector expansion.

\subsection{Finite type-transfer certificate}
The finite shell-ratio theorem is itself reduced to a nonnegative finite-type
transfer calculation. For a component, let $w_n$ be the row vector of
nonnegative type weights at shell $n$ and let $T$ be a nonnegative transfer
matrix. If
\[
  w_{n+1}=w_nT,\qquad \max_i\sum_j T_{ij}\le r<1,
\]
then, with $B_n=\|w_n\|_1$,
\[
  B_{n+1}=\sum_j(w_nT)_j
        =\sum_i w_{n,i}\sum_jT_{ij}
        \le r\sum_iw_{n,i}=rB_n .
\]
The verifier applies this row-sum argument separately to the value, gradient,
and Hessian $\ell^1$ type systems. It then constructs the finite shell-ratio
certificate and invokes the existing ratio-tail and geometric-tail pipeline.
Thus the next Yang--Mills-specific analytic target is a finite Ward-screened
type system whose transfer matrices have row sums below one, with enough
raw-to-coarse and Karp margin downstream.

\subsection{Ward-defect surplus principle}
The type-transfer contraction has a sharper analytic source. For each type
$i$, let $d_i\ge 0$ be the Ward-screening damping gained by the local Ward
cancellation and let $L_{ij}\ge 0$ be the leakage from type $i$ to type $j$
after the screened cavity expansion. Define
\[
  T_{ij}=e^{-d_i}L_{ij}.
\]
If the Ward damping beats branching and leakage in the sense that
\[
  \sum_j e^{-d_i}L_{ij}\le r<1
\]
for every source type $i$, then $T$ is a nonnegative transfer matrix with row
sums at most $r$. The finite type-transfer theorem therefore gives
$B_{n+1}\le rB_n$ for the value, gradient, and Hessian type systems. This
principle is the current analytic breakthrough target: prove a genuine
Yang--Mills Ward-defect surplus estimate, not merely a formal row-sum bound.

\subsection{Interval Ward-defect surplus}
For computer-assisted estimates the Ward-defect surplus is used in interval
form. Suppose each type row has a rigorous lower damping bound $d_i^-$ and
rigorous leakage upper bounds $L_{ij}^+$. Since $e^{-d}$ is decreasing,
\[
  \sum_j e^{-d_i}L_{ij}
  \le e^{-d_i^-}\sum_j L_{ij}^+ .
\]
Thus the endpoint inequality
\[
  e^{-d_i^-}\sum_j L_{ij}^+\le r<1
\]
is enough to produce the Ward-defect surplus certificate. This is the finite
validation target for the analytic Yang--Mills estimates: prove lower bounds
on Ward-screening damping and upper bounds on leakage rows whose endpoint row
sums leave strict surplus.

\subsection{Primitive Ward-screening estimates}
The interval package can itself be produced from primitive estimates with
certified error budgets. If a sector computation gives a damping estimate
\(\widehat d_i\) with error \(\varepsilon_i\), and leakage estimates
\(\widehat L_{ij}\) with errors \(\varepsilon_{ij}\), define
\[
  d_i^-=\widehat d_i-\varepsilon_i,\qquad
  L_{ij}^+=\widehat L_{ij}+\varepsilon_{ij}.
\]
When these outward endpoints are nonnegative, they are rigorous lower damping
and upper leakage bounds. Thus the decisive finite analytic target has become
primitive: prove Yang--Mills Ward-screening estimates whose outward-rounded
endpoints still satisfy
\[
  e^{-(\widehat d_i-\varepsilon_i)}
    \sum_j(\widehat L_{ij}+\varepsilon_{ij})<1.
\]
The remaining problem is not to invent another formal converter, but to prove
these primitive damping and leakage estimates for the actual sector expansion.

\subsection{Relative Ward-screening gain}
A still sharper way to state the analytic input is relative to the unscreened
leakage row. Let
\[
  B_i=\sum_j L_{ij}^+
\]
be the outward upper bound for the baseline leakage mass from type \(i\), and
let \(R_i\) be an upper bound for the residual mass after applying the local
Ward cancellation. If \(0<R_i\le B_i\), define
\[
  d_i=\log(B_i/R_i).
\]
Then
\[
  e^{-d_i}B_i=R_i.
\]
Consequently the whole Ward surplus check is reduced to proving a residual
mass inequality \(R_i\le r<1\) for the value, gradient, and Hessian type
systems. This is the current most primitive finite contract for the analytic
breakthrough: prove that Ward cancellation leaves residual leakage mass below
the contraction threshold.

\subsection{Paired Ward-cancellation residuals}
The residual mass \(R_i\) can be certified by explicit cancelling pairs. Group
the signed Ward contributions in a type row into pairs \((p_k,n_k)\), with
certified error \(\varepsilon_k\), and let \(U_i\) be the unpaired residual
budget. Then
\[
  R_i=\sum_k \bigl(|p_k+n_k|+\varepsilon_k\bigr)+U_i
\]
is a rigorous upper bound on the screened residual. Therefore a finite
pairing whose pair defects and unpaired budget satisfy \(R_i\le r<1\) supplies
the relative Ward-screening certificate. This is the next concrete analytic
construction problem: build the Ward-pairing for the actual Yang--Mills sector
expansion and prove the pair defects are small.

\subsection{Involutive Ward term pairing}
The finite Ward-pairing can be certified by an involution on the signed term
set. An involution partitions the terms into two-cycles and fixed points. Each
two-cycle gives a cancelling pair; the pair error is the sum of the two term
errors. Each fixed point is charged to the unpaired residual by its absolute
term value plus its error. Summing absolute term values and errors by target
type gives the baseline leakage rows. Thus the next primitive construction is
a finite Ward involution on the actual sector terms, with sufficiently small
term errors and few or small fixed points.

\subsection{Fractional Ward cancellation transport}
The failed exact-involution test suggests a strictly weaker finite object.  A
fractional Ward cancellation transport splits each source row of signed sector
terms into positive masses $p_a$ and negative masses $n_b$ and supplies a
nonnegative transport plan $\pi_{ab}$ with row sums at most $p_a$ and column
sums at most $n_b$.  Transported equal mass cancels.  The verifier charges only
the unmatched mass, the certified term errors, any declared residual budget,
and the cost-weighted transport defects
\[
        \sum_{a,b}\pi_{ab}c_{ab}+\epsilon_{ab}.
\]
Thus the screened residual is
\[
 R_i =
 \sum_a\Bigl(p_a-\sum_b\pi_{ab}\Bigr)
 +\sum_b\Bigl(n_b-\sum_a\pi_{ab}\Bigr)
 +\sum_{a,b}(\pi_{ab}c_{ab}+\epsilon_{ab})
 + E_{\rm term}+E_{\rm res}.
\]
If $0<R_i\le B_i$ and $R_i\le r<1$ for each value, gradient, and Hessian type
row, the existing relative Ward-screening certificate applies.  This
generalizes the involutive scheme: an exact two-cycle pairing is the special
case where each transported mass saturates a single positive and negative term
with zero transport defect.  The new object is therefore a conditional route
around the failed exact Ward involution, not a Clay proof.  It replaces a
rigid combinatorial symmetry by a finite linear-program-style cancellation
certificate whose failure or success is checkable.

\subsection{Interpolation-driven fractional Ward transport}
The fractional transport target should not be implemented by solving a
large discrete matching problem after the cluster expansion has already been
enumerated.  The intended Yang--Mills construction is instead an
interpolation-driven fractional Ward transport (IDFWT).  Work in a finite
gauge-fixed cavity chart and let $\mu_s$, $0\le s\le 1$, be a continuous
decoupling interpolation from independent cavities to the physical lattice.
The parameter $s$ may be the edge weight in a Brydges--Kennedy--Rivasseau
forest formula, or an equivalent heat-kernel/Langevin homotopy for the same
finite chart.  On the fully coupled finite chart the Ward identity is exact.
Hence the derivative of the signed Ward defect along $s$ is a conservative
current, up to the explicit boundary and localization errors already charged
in the DLCR ledger.

For a source type row, write the signed infinitesimal Ward defect as
\[
  \dot A_s=\dot A_s^+-\dot A_s^-+\dot E_s .
\]
The IDFWT candidate transport is the integrated positive-to-negative current
\[
  \pi_{ab}=\int_0^1 J_{ab}(s)\,ds,\qquad J_{ab}(s)\ge0,
\]
where $J_{ab}(s)$ is read from the common interpolation event that moves
defect mass from the positive sector label $a$ to the negative sector label
$b$.  In a BKR realization the event is a forest edge contraction; in a
heat-kernel realization it is the corresponding infinitesimal Markov current.
This definition is analytic rather than combinatorial: the transport plan is
built during the decoupling expansion, and its row and column inequalities are
the integrated continuity inequalities for the positive and negative defect
densities.

The finite verifier therefore needs only outward bounds for the current:
\[
  \int_0^1 J_{ab}(s)\,ds\le \pi_{ab}^+,\qquad
  c_{ab}\le c_{ab}^+,\qquad
  \epsilon_{ab}\le \epsilon_{ab}^+ .
\]
The cost $c_{ab}$ is chosen proportional to the interpolated interaction
strength of the edge or heat-kernel step producing the current.  Standard
tree-graph estimates then give a falsifiable sufficient condition
\[
  \sum_{a,b}\pi_{ab}^+c_{ab}^+
  +\sum_{a,b}\epsilon_{ab}^+
  +E_{\rm term}+E_{\rm res}
  +\hbox{unmatched mass}
  \le r_i<1 .
\]
Because the BKR expansion factorizes over forests, this bound is a product of
local tree amplitudes and no longer asks for a one-to-one pairing among all
Feynman/polymer graphs.  The remaining field-theoretic obligation is precise:
construct the finite SU(3) sector expansion, prove the interpolation current
identity in the chosen chart, and certify that exponential tree decay beats
the cost-weighted current and the boundary residuals uniformly in the cutoff.
If that strict inequality survives interval rounding, the existing fractional
Ward transport theorem promotes it to relative Ward screening, then to
primitive Ward estimates, interval Ward-defect surplus, finite type transfer,
shell ratios, geometric cluster tails, and the screened bridge norm.

\subsection{Finite BKR continuity equation}
The finite BKR Ward-current identity can be split into a still smaller
continuity-equation certificate.  On a finite gauge-fixed chart let \(D(s)\)
be the signed Ward defect along a \(C^1\) BKR interpolation.  The analytic
identity begins with the fundamental theorem of calculus,
\[
  D(1)-D(0)=\int_0^1D'(s)\,ds .
\]
The sectorized certificate writes \(D'(s)\) as a finite signed sum of forest
current densities plus declared gauge, boundary, truncation, and balance
errors.  For each source type row the verifier computes
\[
  \Delta_{\rm end}
  =\sum_a p_a-\sum_b n_b,\qquad
  I_{\rm forest}
  =\sum_F \sigma_F\int_0^1J_F(s)\,ds ,
\]
where \(\sigma_F=\pm1\) records the orientation of the forest current in the
signed continuity equation.  The row is accepted if
\[
  |\Delta_{\rm end}-I_{\rm forest}|
  \le
  E_{\rm gauge}+E_{\rm bdry}+E_{\rm trunc}+E_{\rm bal}.
\]
This proves only the finite signed balance.  It does not yet prove that all
oriented currents have the positive-to-negative form needed for IDFWT, nor
does it prove the tree-decay constants.  But it isolates the first analytic
gate: finite-chart \(C^1\) BKR differentiability plus a sectorized current
decomposition with explicit balance error.

\subsection{Six gate breakdown}
The continuity package now splits into six named gates.  The verifier treats
them as separate obligations:
\begin{enumerate}
\item finite \(C^1\) interpolation;
\item sectorized current decomposition;
\item endpoint-current balance;
\item usable-current orientation;
\item tree-decay cost estimate;
\item cutoff-stable residual margin.
\end{enumerate}
This is the finer roadmap the project needed.  Each gate has its own field in
the package, so the proof no longer hides a single large interpolation claim
behind a monolithic theorem.  Instead, the remaining SU(3) work is now
distributed into six smaller obligations that can be attacked one by one.

\subsection{Genuine SU(3) continuity package roadmap}
The six gates can be made still more operational.  A genuine
\(\mathrm{SU}(3)\) finite BKR continuity package must provide the following
finite subcertificates.

For the finite \(C^1\) interpolation gate, one must give a tree-gauge chart
atlas on \(SU(3)^m/\operatorname{Ad}\), a BKR link-weight interpolation
\(s\mapsto S_s\) with the common action hash, uniform \(C^1\) domination for
the action density and Ward observable, and a finite quadrature or interval
enclosure rule for \(\int_0^1D'(s)\,ds\).  The verifier-facing acceptance test
is that the package has
\(\texttt{differentiability\_status}=\texttt{FINITE\_CHART\_C1}\) and finite
derivative bounds on the declared parameter box.

For the sectorized current decomposition gate, one must give a finite sector
partition of the plaquette and edge words, a signed endpoint Ward-term table
with target types, a forest-current event table for each derivative term, and
a residual budget for discarded or merged terms.  The acceptance test is that
the sector decomposition is declared finite and every continuity row has
positive terms, negative terms, and finite forest-current events.

For the endpoint-current balance gate, one must give the endpoint defect
\(D(1)-D(0)\), the integrated signed forest current, and the gauge-fixing,
boundary, truncation, and balance error budgets.  The finite inequality is
\[
  \bigl|\Delta_{\rm end}-I_{\rm forest}\bigr|
  \le
  E_{\rm gauge}+E_{\rm bdry}+E_{\rm trunc}+E_{\rm bal}.
\]

For the usable-current orientation gate, one must give positive endpoint
capacities, negative endpoint capacities, an orientation sign for each
forest-current event, and row/column capacity inequalities for the IDFWT
transport plan.  The acceptance test is that every consumed current is
nonnegative after orientation and moves mass from a positive Ward-defect label
to a negative Ward-defect label without exceeding the two capacities.

For the tree-decay cost gate, one must give for each event an interaction
strength upper bound \(g_F\), a tree amplitude \(A_F\), a decay ratio
\(\rho_F<1\), a tail start \(n_F\), and a quadrature error.  The cost charged
to the transport event is then bounded by
\[
  g_F\,{A_F\rho_F^{n_F}\over 1-\rho_F}.
\]

For the cutoff-stable residual-margin gate, one must give unmatched-mass
bounds, the cost-weighted transport defect sum, outward interval-rounding
budgets, and cutoff/volume variation budgets with strict margin.  The final
acceptance test is the screened contraction inequality
\[
  R_i\le r_i<1
\]
for every value, gradient, and Hessian type row after all errors are charged.

Thus the next genuine theorem is no longer a vague continuity assertion.  It
is a finite construction problem: populate these six subcertificate families
with a matching \(\mathrm{SU}(3)\) action hash and strict margins.  If that is
done, the existing finite BKR continuity converter emits the finite-chart BKR
Ward-current identity, then IDFWT, fractional Ward transport, and the
downstream screened DLCR certificates.  This paragraph is still only a
roadmap reduction; it does not prove the \(\mathrm{SU}(3)\) chart atlas, the
\(C^1\) estimates, the signed current identity, the tree-decay constants, or
the cutoff-stable residual margin.

\subsection{Micro-roadmap for the six SU(3) gates}
The subcertificate list can be turned into a still finer workbench.  The first
gate begins by fixing a packet graph and maximal tree gauge, writing the
non-tree holonomies as \(SU(3)^m\) coordinates modulo residual conjugation,
defining the BKR edge-weight action \(S_s\) on a compact parameter box,
differentiating the finite-dimensional integral under the integral sign, and
recording derivative majorants for the value, gradient, and Hessian rows.  Its
first decisive test is one nontrivial plaquette-edge chart with a certified
\(C^1\) derivative enclosure and matching action hash.

The second gate expands each differentiated action-observable product into
finite word sectors, assigns every signed endpoint contribution to a source
and target type, names the BKR forest event that generated it, merges
symmetry-equivalent terms only with an explicit residual budget, and proves
that no untyped finite term remains outside the residual.  Its first decisive
test is a complete signed term table for the same raw edge used in the first
gate.

The third gate computes the endpoint signed Ward defect, integrates the signed
forest-current table over the BKR parameter, separates gauge-fixing, boundary,
truncation, and quadrature errors, proves the row balance inequality by
outward interval arithmetic, and stores both the balance defect and the
allowed defect.  Its first decisive test is that the largest row balance
defect is below the declared continuity-error budget.

The fourth gate splits the signed table into positive and negative endpoint
capacities, chooses an orientation for each forest event, discards or charges
currents that cannot be oriented from positive to negative mass, and checks
the row and column capacity inequalities.  Its first decisive test is that the
usable oriented current transports a positive fraction of every dominant
positive defect without exceeding the negative capacities.

The fifth gate bounds the local interpolated interaction strength, proves a
tree-amplitude bound, proves a geometric decay ratio strictly below one,
chooses a tail start order compatible with the truncation, and rounds the cost
envelope outward before passing to IDFWT.  Its first decisive test is that the
cost-weighted oriented current is below the available Ward-screening surplus
on the test raw edge.

The sixth gate sums unmatched endpoint mass, cost-weighted transport defects,
continuity errors, interval-rounding reserves, rationalization reserves, and
cutoff/volume variation reserves.  It then checks the residual contraction and
the downstream Karp margin simultaneously.  Its first decisive test is that
the value, gradient, and Hessian rows all satisfy \(R_i\le r_i<1\) after all
reserves and that the emitted coarse graph still has positive Karp margin.

This micro-roadmap gives a natural order for the next proof attempts:
\begin{align*}
  C^1\hbox{ chart}
  &\to \hbox{typed sectors}
  \to \hbox{signed balance}\\
  &\to \hbox{usable orientation}
  \to \hbox{tree cost}
  \to \hbox{residual and Karp margin}.
\end{align*}
Failure at any arrow is informative: it identifies the first analytic
obstruction before the project spends effort on downstream certificate
machinery.

\subsection{Artifact-level SU(3) continuity roadmap}
The micro-roadmap can be refined once more into named mini-artifacts.  Each
mini-artifact has four fields: the lemma to prove, the finite data it emits,
the verifier fields it populates, and the failure mode that would stop the
route.  The intended order is:
\begingroup\small
\begin{enumerate}
\item \texttt{su3\_tree\_gauge\_chart.v1}: construct the finite
  \(SU(3)^m/\operatorname{Ad}\) chart and emit coordinates and boxes;
\item \texttt{su3\_bkr\_action\_path.v1}: prove the edge-weight path
  \(S_s\) is \(C^1\) with the common action hash;
\item \texttt{su3\_ward\_derivative\_enclosure.v1}: enclose \(D'(s)\) for
  value, gradient, and Hessian rows;
\item \texttt{su3\_word\_sector\_partition.v1}: split differentiated words
  into finite source and target sectors;
\item \texttt{su3\_signed\_endpoint\_terms.v1}: certify signs, target types,
  and error radii for endpoint Ward terms;
\item \texttt{su3\_forest\_current\_table.v1}: attach each derivative term to
  a BKR forest-current event;
\item \texttt{su3\_endpoint\_defect\_interval.v1}: evaluate \(D(1)-D(0)\)
  with outward intervals;
\item \texttt{su3\_integrated\_current\_interval.v1}: integrate the signed
  forest current with quadrature error;
\item \texttt{su3\_balance\_error\_budget.v1}: prove the rowwise balance
  inequality;
\item \texttt{su3\_endpoint\_capacity\_table.v1}: bound positive and negative
  endpoint capacities;
\item \texttt{su3\_oriented\_current\_selection.v1}: select the usable
  positive-to-negative subcurrent and charge the rest;
\item \texttt{su3\_idfwt\_capacity\_check.v1}: verify the IDFWT row and
  column capacity inequalities;
\item \texttt{su3\_link\_coupling\_bound.v1}: bound each active interpolated
  link strength;
\item \texttt{su3\_tree\_tail\_envelope.v1}: prove \(A_F\rho_F^n/(1-\rho_F)\)
  with \(\rho_F<1\);
\item \texttt{su3\_transport\_cost\_rounding.v1}: round the cost envelope
  outward while preserving surplus;
\item \texttt{su3\_unmatched\_mass\_accounting.v1}: charge every
  untransported endpoint mass by type row;
\item \texttt{su3\_uniform\_reserve\_budget.v1}: bound rounding,
  rationalization, cutoff, and volume reserves;
\item \texttt{su3\_residual\_karp\_margin.v1}: prove \(R_i\le r_i<1\) and a
  positive coarse Karp margin together.
\end{enumerate}
\endgroup
This is now the smallest useful roadmap unit in the current repository.  A
failure of \(\texttt{su3\_tree\_tail\_envelope.v1}\), for example, means the
tree branching has beaten decay; a failure of
\(\texttt{su3\_residual\_karp\_margin.v1}\) means local residual contraction
did not survive the global coarse-graph drift test.  Neither failure would be
ambiguous, and neither would require re-auditing the whole DLCR chain.

\subsection{Atomic-obligation roadmap}
Each of the eighteen mini-artifacts is now split into three atomic
obligations.  They are, respectively, data normalization, analytic estimate,
and verifier handoff.  This gives fifty-four atomic tasks.  For example, the
tree-gauge chart artifact breaks into:
\begin{enumerate}
\item choose the finite packet graph, root set, and maximal tree;
\item prove every packet loop word is represented by non-tree \(SU(3)\)
  coordinates modulo residual conjugation;
\item emit coordinate boxes and raw-edge support with the matching action hash.
\end{enumerate}
At the other end,
\(\texttt{su3\_residual\_karp\_margin.v1}\) breaks into:
\begin{enumerate}
\item verify \(R_i\le r_i<1\) for value rows;
\item verify the same residual contraction for gradient and Hessian rows;
\item verify that the emitted coarse transition graph has positive Karp
  margin.
\end{enumerate}
The intermediate artifacts follow the same pattern.  The BKR action-path
artifact must declare active links, differentiate the Wilson action path, and
check the two endpoint actions.  The Ward-derivative enclosure must write
\(D'(s)\), dominate the differentiated gradient/Hessian representatives, and
produce quadrature intervals.  The sector artifacts must define the type
alphabet, assign every differentiated word, certify endpoint signs, and charge
ambiguous or unassigned terms.  The current-balance artifacts must compute
endpoint and integrated-current intervals and then prove
\(\text{balance\_defect}\le\text{allowed\_defect}\).  The orientation artifacts
must record capacities, choose a nonnegative positive-to-negative subcurrent,
and verify IDFWT row and column inequalities.  The cost artifacts must bound
active link strength, prove \(A_F,\rho_F,n_F\) with \(\rho_F<1\), and show
that rounded cost defects remain below surplus.  The residual artifacts must
account for unmatched mass, add all reserve budgets, and close contraction
together with Karp negativity.

This atomic layer is useful because it identifies the next smallest proof
object: before trying to prove any Ward transport estimate, the project must
first produce the first three atomic obligations of
\(\texttt{su3\_tree\_gauge\_chart.v1}\).  If even that chart artifact cannot
be made genuine on a nontrivial raw edge, the IDFWT route has no honest
finite-chart input.

\subsection{Leaf plan for the first chart artifact}
The immediate next artifact,
\(\texttt{su3\_tree\_gauge\_chart.v1}\), can itself be split into nine leaf
certificates:
\begin{enumerate}
\item finite graph inventory: list vertices, oriented edges, packet paths, and
  the raw edge under test;
\item root and tree choice: choose roots and a maximal tree or forest, with
  \(|T|=|V|-\)components and no cycles;
\item tree-gauge section: construct the representative setting tree holonomies
  to the identity;
\item non-tree coordinate map: identify the remaining holonomies as
  \(SU(3)^m\) coordinates with \(m=|E|-|V|+\)components;
\item loop-word reconstruction: rewrite every packet loop and tested raw edge
  as a word in the non-tree coordinates;
\item residual adjoint action: record root conjugation on \(SU(3)^m\) and
  trace-word invariance;
\item compact chart domain: declare the compact domain and the finite interval
  boxes used by later enclosures;
\item raw-edge support binding: bind the source and target raw states of the
  test edge to the chart partition;
\item action-hash binding: attach the common Yang--Mills action hash and
  chart-local action metadata.
\end{enumerate}
This leaf plan is deliberately elementary.  It is the first point at which the
roadmap can stop being only a dependency graph and become a genuine finite
artifact.  The next mathematical move is to instantiate these nine leaves for
one nontrivial raw edge; only then does it make sense to attempt the BKR
action path and Ward-derivative enclosure artifacts.

\subsection{Field plan for the first leaf}
The first leaf, finite graph inventory, can be refined into eight fields:
\[
\begin{gathered}
\text{vertices},\quad \text{oriented edges},\quad \text{packet paths},\quad
\text{packet loops},\\
\text{raw edge under test},\quad \text{plaquette words},\quad
\text{component labels},\quad \text{inventory hash}.
\end{gathered}
\]
The acceptance checks are intentionally elementary.  The vertex set must be
nonempty, every edge endpoint must be declared, every packet path must use
declared oriented edge labels, every packet loop and plaquette word must be
closed, the raw edge under test must be supported by the packet graph, and the
component labels must agree with graph connectivity.  Finally, the inventory
hash must be canonical under the ordering conventions used by the verifier.
This first leaf has four immediate failure modes: a packet path names an
undeclared edge, the raw edge support lies outside the packet graph, a
plaquette boundary is not closed, or the component labels disagree with graph
connectivity.  If the field plan passes, it hands the vertices, oriented
edges, and component labels to the next leaf, where the maximal-tree check is
\(|T|=|V|-\)components.

\subsection{Finite-chart BKR Ward-current identity}
The next roadmap layer feeds IDFWT from the continuity equation.  A finite
BKR-current package records, for each finite chart and source type row, the
positive and negative Ward endpoint masses together with BKR forest-current
events.  After the signed balance above has been verified, the package records
the finite-chart current balance
\[
   A_1^+-A_1^- - (A_0^+-A_0^-)
   =
   \sum_F\int_0^1 J_F(s)\,ds
   + E_{\rm gauge}+E_{\rm bdry}+E_{\rm trunc},
\]
after sectorization and gauge fixing.  Each nonnegative forest current
\(J_F\) names a positive endpoint \(a\), a negative endpoint \(b\), and an
interpolated link or block coupling.  The finite verifier converts the event
to an IDFWT transport mass
\[
  \pi_{ab}^F=\int_0^1J_F(s)\,ds
\]
and records the continuity error
\[
  E_{\rm cont}=E_{\rm gauge}+E_{\rm bdry}+E_{\rm trunc}.
\]
If the event has link-coupling upper bound \(g_F\), tree amplitude \(A_F\),
decay ratio \(\rho_F<1\), and tail start \(n_F\), its transport cost is
bounded by
\[
  c_F \le g_F\,{A_F\rho_F^{n_F}\over 1-\rho_F}.
\]
Thus a proved finite BKR Ward-current package canonically emits the IDFWT
package; the already stated IDFWT verifier then checks capacities, costs,
unmatched mass, and residual contraction.  The remaining genuine
Yang--Mills obligation has therefore been sharpened again: prove finite
\(C^1\) BKR differentiability, construct the sectorized signed current
decomposition, orient the useful currents from positive to negative Ward
mass, and certify cutoff-stable tree-decay constants.

\subsection{Roadmap audit}
The current proof state separates two facts that must not be conflated. The
final bridge registry is discharged as a registry matter by explicit
algebraic/GNS and spectral/generator witness artifacts, but the finite-chart
Yang--Mills analytic input is not constructed. A direct continuum
OS-identification route leaves the compactness, localization, and spectral
identification problems in essentially their original form. The Ward-screened
DLCR route is therefore the selected conditional framework: its recent
converter and stability theorems reduce the remaining field-theoretic input to
a finite Ward-projected interval jet atlas, raw-to-coarse domination with
positive slack, a negative coarse Karp cycle-mean certificate, and uniform
cutoff/volume stability. This is a promising roadmap, not a proof of the Clay
conjecture. The next decisive theorem must construct a genuine Yang--Mills
BKR continuity package whose emitted Ward-current and IDFWT data have screened
residual $R_i\le r_i<1$ and whose interval-rounded outputs leave enough tail
slack and Karp margin to feed this finite chain. The six named subgates above
are the finer checklist for that theorem.

\subsection{Continue-or-pivot decision}
The current probability of completing the Clay proof from the present state
should be regarded as below one half. The reason is not lack of progress:
recent reductions have made the next obstruction finite and falsifiable. The
reason is that the genuine Yang--Mills finite \(C^1\) BKR continuity equation,
sector term estimates, tree-decay constants, negative Karp certificate, and
uniform cutoff/volume stability have not been constructed. The correct
decision after the exact-involution failure is therefore a narrow pivot: stop
searching for a raw exact Ward involution and attempt one nontrivial
finite-chart BKR continuity package whose emitted IDFWT instance has certified
residual contraction and positive downstream margins. If that
continuous transport test also fails, the project should remain conditional
verifier infrastructure rather than an open-ended Clay proof claim.

\subsection{Decisive-test outcome}
The decisive raw-edge Ward-involution attempt fails in the present repository.
There is no proved genuine Yang--Mills artifact with schema
\texttt{yangmills.dlcr.involutive\_ward\_term\_pairing.v1}, and no available
sector term list, target-type assignment, certified term-error table, or Ward
involution from which such an artifact can be constructed. The Yang--Mills
evidence files that could feed this object remain explicitly conditional.
Therefore this route should no longer be presented as an active likely Clay
proof path. The useful output is conditional verifier infrastructure and a
record of finite reductions. The next analytic search target is not the
exact involutive Ward-pairing artifact; it is a genuine finite-chart BKR
continuity package whose emitted Ward-current and IDFWT data produce the
fractional Ward transport artifact.

\subsection{Conditional implication}
If the exact shadow coupling, reflection-time intertwiner, and
no-current-leakage lemmas are proved for the same action specification and the
same OS/GNS representation, then the convergence-side Schwinger realization on
$W_N$ and the convergence-side action realization on $F_N$ follow. The
conditional implication is recorded in the DLCR proof artifact. The repository
does not currently prove those remaining analytic DLCR lemmas. The audited status remains theorem boundary.
\endinput
